% !TeX root = ../main.tex

\documentclass[../main.tex]{subfiles}
\begin{document}

\chapter{TÍNH TOÁN CÂN BẰNG VẬT CHẤT VÀ NĂNG LƯỢNG TOÀN HỆ THỐNG}
\label{chuong3}

%%=============================================================
\section{Sơ đồ công nghệ và thông số cơ sở}
%%=============================================================

%%=============================================================

\subsection{Sơ đồ dòng chảy công nghệ tổng thể (PFD)}
\label{sec:so-do-pfd}

Sơ đồ dòng chảy công nghệ (Process Flow Diagram -- PFD) tổng thể thể hiện toàn bộ chuỗi xử lý từ CTRSH đầu vào đến sản phẩm cuối cùng (điện, dầu DO, than carbon, khí thải và nước thải), được xây dựng dựa trên các kết quả cân bằng vật chất và năng lượng đã tính toán trong các mục trên.

\begin{figure}[htbp]
  \centering
  \begin{adjustbox}{max width=1.1\textwidth, center}
  \begin{tikzpicture}[
      node distance=0.5cm and 0.4cm,
      >=Stealth,
      flow/.style={draw, rounded corners=3pt, text width=3.6cm, minimum height=1.0cm, font=\small, align=center, inner sep=4pt},
      storage/.style={draw, diamond, aspect=2, minimum width=2.4cm, minimum height=1.4cm, font=\small, align=center},
      stream/.style={draw, rounded corners=2pt, minimum width=2.2cm, minimum height=0.7cm, font=\tiny, align=center, fill=blue!8},
      arrow/.style={->, >=Stealth, very thick},
      arr/.style={->, >=Stealth, thick, dashed},
      unit/.style={rectangle, draw, rounded corners=2pt, minimum width=1.8cm, minimum height=0.6cm, font=\tiny, align=center},
      legend/.style={rectangle, draw, font=\footnotesize, align=left, fill=white, inner sep=4pt}
  ]
  
  \node[legend, anchor=south west] (leg) at (-1.8, 1.5) {
      \begin{tabular}{l}
      \textbf{Chú thích:}\\
      \textcolor{green!70!black}{\rule[0.5ex]{3mm}{2mm} Dòng vật chất}\\
      \textcolor{blue!70!black}{\rule[0.5ex]{3mm}{2mm} Dòng năng lượng}\\
      \end{tabular}
  };
  
  \node[flow, fill=green!10, draw=green!60!black, anchor=south east, text width=6cm] (sanpham) at (15.3, 1.5) {
      \begin{tabular}{l}
      \textbf{Sản phẩm chính:}\\
      $\bullet$ Điện năng: $\approx$921 kW (ròng)\\
      $\bullet$ Dầu DO: $\approx$10,93 t/ngày\\
      $\bullet$ Than carbon: 11,82 t/ngày
      \end{tabular}
  };
  
  % LỚP 1 (Y = 0)
  \node[flow, fill=green!15, draw=green!70!black] (ctrs) at (0, 0) {CTRSH Hội An\\(120 tấn/ngày)};
  \node[flow, fill=blue!10, draw=blue!70!black] (pit) at (4.5, 0) {Pit tiếp nhận\\\& Hút mùi};
  \node[flow, fill=orange!15, draw=orange!70!black] (phanloai) at (9, 0) {Phân loại cơ học\\Tách từ, sàng};
  
  \draw[arrow, green!70!black] (ctrs.east) -- node[midway, above] {} (pit.west);
  \draw[arrow, green!70!black] (pit.east) -- node[midway, above] {} (phanloai.west);
  
  % LỚP 2 (Y = -2.5) TÁCH DÒNG
  \node[flow, fill=red!10, draw=red!60!black] (rac-chat) at (4.5, -2.5) {Chất thải không cháy\\(16,08 t/ngày)};
  \node[flow, fill=green!15, draw=green!70!black] (huucokhong) at (9, -2.5) {Hỗn hợp cháy được\\($\sim$50,3 t/ngày)};
  \node[flow, fill=purple!10, draw=purple!70!black] (nhua-giay) at (13.5, -2.5) {RDF cháy được sau tách lọc\\($\sim$35,5 t/ngày)};
  
  \draw[thick, green!70!black] (phanloai.south) -- (9, -1.25);
  \draw[thick, green!70!black] (4.5, -1.25) -- (13.5, -1.25);
  \draw[arrow, green!70!black] (4.5, -1.25) -- (rac-chat.north);
  \draw[arrow, green!70!black] (9, -1.25) -- (huucokhong.north);
  \draw[arrow, green!70!black] (13.5, -1.25) -- (nhua-giay.north);
  
  % LỚP 3 (Y = -5.0) TIỀN XỬ LÝ
  \node[flow, fill=yellow!15, draw=yellow!70!black] (bien_phap) at (13.5, -5) {Tiền xử lý RDF\\Cắt, sấy, tạo viên};
  \node[flow, fill=cyan!10, draw=cyan!70!black] (nuoc-ep) at (9, -5) {Hơi nước sấy\\(7,88 t/ngày)};
  
  \draw[arrow, green!70!black] (nhua-giay.south) -- (bien_phap.north);
  \draw[arrow, green!70!black] (bien_phap.west) -- (nuoc-ep.east);
  
  % LỚP 4 (Y = -7.5)
  \node[flow, fill=yellow!20, draw=yellow!70!black] (rdf-say) at (13.5, -7.5) {RDF sau sấy\\(39,41 t/ngày\\Độ ẩm $<15\%$);};
  \node[flow, fill=cyan!15, draw=cyan!70!black] (xl-nuoc) at (9, -7.5) {Hệ thống XLNT\\MBBR -- Sinh học};
  
  \draw[arrow, green!70!black] (bien_phap.south) -- node[midway, right] {39,41 t/n} (rdf-say.north);
  \draw[arrow, green!70!black] (nuoc-ep.south) -- node[midway, right, font=\footnotesize] {Hơi nước} (xl-nuoc.north);
  
  % LỚP 6 (Y = -12.5)
  \node[flow, fill=teal!15, draw=teal!70!black] (khi-khong-ngung) at (9, -12.5) {Khí không ngưng\\13,79 t/ngày};
  \node[flow, fill=gray!10, draw=gray!60!black] (xl-khi) at (13.5, -12.5) {Hệ thống XL khí thải\\ESP -- Tháp hấp thụ -- SNCR};
  \node[flow, fill=blue!20, draw=blue!70!black] (chung-cat) at (4.5, -12.5) {Chưng cất\\\& Tinh chế dầu\\H$_2$SO$_4$--NaOH--Đất sét};
  \node[flow, fill=green!20, draw=green!70!black] (dau-do) at (0, -12.5) {Dầu nhiệt phân\\tinh chế\\11,03 t/ngày};

  % LỚP 5 (Y = -10.0) LÒ NHIỆT PHÂN
  \node[flow, fill=orange!20, draw=orange!80!black] (lo-np) at (9, -10) {Lò nhiệt phân quay\\400--500$^\circ$C ($\times$5 mô-đun)};
  \node[flow, fill=blue!15, draw=blue!70!black] (dau-tho) at (4.5, -10) {Dầu thô\\13,79 t/ngày};
  \node[flow, fill=gray!20, draw=gray!70!black] (than-c) at (13.5, -10) {Than carbon\\11,82 t/ngày};

  \draw[arrow, green!70!black] (rdf-say.south) -- (13.5, -8.75) -- (9, -8.75) -- (lo-np.north);

  \draw[arrow, green!70!black] (lo-np.west) -- (dau-tho.east);
  \draw[arrow, green!70!black] (lo-np.east) -- (than-c.west);

  \draw[arrow, green!70!black] (than-c.south) -- (xl-khi.north);

  \draw[arrow, green!70!black] (dau-tho.south) -- (chung-cat.north);
  \draw[arrow, green!70!black] (chung-cat.west) -- (dau-do.east);

  \draw[arrow, green!70!black] ([xshift=0.5cm]lo-np.south) -- ([xshift=0.5cm]khi-khong-ngung.north);
  \draw[arrow, blue!70!black, dashed] ([xshift=-0.5cm]khi-khong-ngung.north) -- ([xshift=-0.5cm]lo-np.south) node[midway, left, font=\scriptsize] {Khí hồi lưu};
  
  % LỚP 7 (Y = -15.0) PHÁT ĐIỆN
  \node[flow, fill=gray!5, draw=gray!40!black] (khi-sau-xl) at (13.5, -15) {Khí thải sau xử lý\\\emph{Đạt QCVN 19:2024}};
  \node[flow, fill=red!15, draw=red!70!black] (diesel) at (0, -15) {Động cơ diesel\\6 $\times$ 550 kVA};
  \node[flow, fill=yellow!20, draw=yellow!80!black] (dien) at (4.5, -15) {\textbf{Điện năng}\\1,84 MW (ròng)};
  
  \draw[arrow, green!70!black] (xl-khi.south) -- (khi-sau-xl.north);
  \draw[arrow, green!70!black] (dau-do.south) -- node[midway, right, font=\footnotesize] {Dầu DO} (diesel.north);
  \draw[arrow, green!70!black] (diesel.east) -- (dien.west);
  
  % LỚP 8 (Y = -17.5)
  \node[flow, fill=gray!10, draw=gray!60!black] (khi-diesel) at (0, -17.5) {Khí thải diesel\\NO$_x$, CO, bụi PM};
  \draw[arrow, green!70!black] (diesel.south) -- (khi-diesel.north);

  % === KẾT NỐI CÁC LỚP ===
  \draw[arrow, green!70!black] (xl-khi.south) -- (khi-sau-xl.north);
  \end{tikzpicture}
  \end{adjustbox}
  \caption{Sơ đồ dòng chảy công nghệ (PFD) tổng thể nhà máy nhiệt phân -- chưng cất -- phát điện diesel Hội An}
  \label{fig:pfd-tong-the}
\end{figure}

Sơ đồ PFD trên thể hiện chuỗi công nghệ liên hoàn: CTRSH $\to$ phân loại cơ học $\to$ tiền xử lý RDF (cắt, sấy, tạo viên) $\to$ lò nhiệt phân quay $\to$ chưng cất dầu $\to$ động cơ diesel phát điện. Các dòng phụ gồm chất thải không cháy (16,08 t/ngày), hơi nước sấy (7,88 t/ngày) và khí thải từ lò NP qua hệ xử lý khí thải nhiều cấp (ESP -- tháp hấp thụ -- SNCR). Khí không ngưng (13,79 t/ngày) được hồi lưu về buồng đốt của lò nhiệt phân, đảm bảo hệ thống có khả năng tự cân bằng năng lượng ở chế độ vận hành ổn định.


\subsection{Cơ sở và các giả thiết tính toán}
%%=============================================================

\subsubsection{Nguyên tắc tính toán cân bằng}

Toàn bộ tính toán cân bằng vật chất được xây dựng trên nguyên tắc bảo toàn khối lượng. Theo nguyên tắc này, tổng khối lượng các dòng đầu vào bằng tổng khối lượng các dòng đầu ra tại mỗi ranh giới hệ thống được xét. Phương trình cân bằng được biểu diễn như sau:
\begin{equation}
  \sum m_{\text{vào}} = \sum m_{\text{ra}} + \Delta m_{\text{tích lũy}}
\end{equation}

Trong điều kiện vận hành ổn định, lượng tích lũy trong hệ thống $\Delta m_{\text{tích lũy}} = 0$, tức là khối lượng ra bằng khối lượng vào. Áp dụng cho từng thiết bị hay khối công đoạn, phương trình bảo toàn khối lượng cụ thể hóa thành:
\begin{equation}
  m_{\text{CTRSH}} = m_{\text{RDF}} + m_{\text{không cháy}} + m_{\text{nước thải, ép}} + m_{\text{hơi nước, bức xạ}} + m_{\text{hơi nước, sấy}} + m_{\text{hao hụt}}
\end{equation}
Phương trình trên là cơ sở cho toàn bộ tính toán cân bằng trong Mục~\ref{sec:can-bang-tong-the}. Khi áp dụng cho khối nhiệt phân, phương trình tương tự được viết lại:
\begin{equation}
  m_{\text{RDF}} = m_{\text{dầu thô}} + m_{\text{khí không ngưng}} + m_{\text{than carbon}}
\end{equation}

Nguyên tắc này đảm bảo mọi dòng vật chất đều được hạch toán đầy đủ, không bỏ sót thành phần nào trong quá trình tính toán. Trong vận hành thực tế, các hao hụt không thể tránh khỏi như bụi, xỉ và vết ngưng bám thành thiết bị cần được ghi nhận riêng và quản lý theo quy định về chất thải nguy hại khi hàm lượng kim loại nặng vượt ngưỡng cho phép.

Nhiệt trị của RDF phụ thuộc mạnh vào thành phần nguyên tố và độ ẩm làm việc. Công thức Mendeleev được sử dụng để ước tính nhiệt trị thấp (LHV) từ thành phần nguyên tố theo phân tích cơ bản \cite{quyen2025giaotrinh}:
\begin{equation}
  Q^r_p = 81 \cdot C^r + 246 \cdot H^r - 26(O^r - S^r) - 6 \cdot W^r \quad [\text{kcal/kg}]
  \label{eq:mendeleev}
\end{equation}
trong đó $C^r,\, H^r,\, O^r,\, S^r,\, W^r$ lần lượt là phần trăm carbon, hydro, oxy, lưu huỳnh và ẩm trong mẫu rác ở trạng thái làm việc. Khi độ ẩm thực tế $W^r$ thay đổi, nhiệt trị quy đổi về trạng thái khô tuyệt đối rồi tính ngược lại về trạng thái làm việc mới theo hệ thức:
\begin{equation}
  Q^r_{p,\,W_2} = Q^r_{p,\,W_1} \cdot \frac{100 - W_2}{100 - W_1} - 6 \cdot (W_2 - W_1)
\end{equation}

Đối với RDF từ CTRSH đô thị ở độ ẩm làm việc 8--15\%, các tài liệu tham khảo cho LHV trong khoảng 4.000--5.500 kcal/kg (16.700--23.000 kJ/kg), với giá trị trung bình thường nằm trong dải 4.300--4.800 kcal/kg \cite{Arena2012,Sharuddin2016}. Trong tính toán tổng quát, LHV của RDF được biểu diễn bằng ký hiệu $Q_{\text{RDF}}^r$ và sẽ được thay bằng số liệu thực nghiệm của từng dự án cụ thể khi áp dụng.

\subsubsection{Các giả thiết tính toán}

Để tính toán sơ bộ, các giả thiết sau được áp dụng nhất quán trong toàn chương:

\begin{enumerate}
  \item Nhà máy vận hành 8.000 giờ/năm, tương đương khoảng 333 ngày vận hành thực tế.
  \item Các lò nhiệt phân vận hành theo mẻ. Tuy nhiên, vì hệ thống gồm nhiều mô-đun lò song song nên khí không ngưng từ lò đang hoạt động có thể cấp nhiệt cho lò đang trong giai đoạn khởi động hoặc làm nóng sấy. Do đó, toàn bộ hệ thống nhiệt phân có thể xem là hoạt động liên tục ở trạng thái ổn định.
  \item Thành phần và độ ẩm của CTRSH đầu vào là ổn định theo giá trị trung bình; biến động theo mùa được xem xét trong phân tích độ nhạy tại Mục~\ref{sec:kiem-tra-hop-ly}.
  \item Tỷ lệ sản phẩm nhiệt phân (dầu -- khí -- than) được tham chiếu từ các kết quả thực nghiệm trên RDF gốc nhựa/nylon \cite{Sharuddin2016,Miandad2016} và lấy theo ký hiệu $y_{\text{dầu}}, y_{\text{khí}}, y_{\text{than}}$ với dải điển hình lần lượt 35--45\%, 25--35\% và 20--30\% khối lượng RDF cấp lò.
  \item Hệ thống vận hành ở trạng thái ổn định, các hiện tượng quá độ khi khởi động và tắt lò không được tính vào cân bằng.
  \item Các tổn thất nhiệt qua thành lò và đường ống được bao hàm trong hiệu suất tổng hợp $\eta \in [0{,}28; 0{,}35]$ mà không tính riêng lẻ từng thành phần.
  \item Toàn bộ số liệu tính toán (nhiệt trị, thành phần nguyên tố, tỷ lệ sản phẩm, hiệu suất, hệ số an toàn) được thống nhất theo các giá trị nêu tại Bảng~\ref{tab:thong-so-dau-vao}.
\end{enumerate}

Các số liệu đầu vào của bài toán cân bằng vật chất mang sai số nhất định. Mức độ không chắc chắn của các thông số chính được đánh giá như sau:

\begin{itemize}
  \item \textbf{Độ ẩm CTRSH}: dao động theo mùa trong khoảng $\pm 8\%$ so với giá trị trung bình. Vào mùa mưa, độ ẩm có thể đạt 60--65\%, làm tăng tải nhiệt của công đoạn sấy và giảm sản lượng RDF.
  \item \textbf{Nhiệt trị RDF}: sai số đo lường từ phòng kiểm nghiệm khoảng $\pm 3\%$; sai số do biến động thành phần nguyên liệu theo mùa ước tính thêm $\pm 5\%$; tổng sai số hợp lý $\pm 8\%$.
  \item \textbf{Tỷ lệ sản phẩm nhiệt phân}: độ không chắc chắn lớn nhất trong toàn bộ tính toán, ước tính $\pm 10\%$ so với giá trị thiết kế, phụ thuộc vào nhiệt độ lò, thời gian lưu và biến động thành phần nhựa trong RDF.
  \item \textbf{Hiệu suất chuyển hóa $\eta$}: sai số $\pm 15\%$ tùy thuộc vào chất lượng dầu sau tinh chế và đặc tính vận hành của động cơ diesel ở tải trọng thực.
\end{itemize}

Trong thiết kế sơ bộ, phương pháp xử lý độ không chắc chắn được áp dụng là sử dụng hệ số an toàn (hệ số không ổn định $K = 1{,}25$) để quy đổi công suất lý thuyết về công suất vận hành thực tế, kết hợp với phân tích độ nhạy tại Mục~\ref{sec:kiem-tra-hop-ly}. Phương pháp này cho phép định lượng mức thay đổi của kết quả đầu ra theo biến động của thông số đầu vào, qua đó xác định các thông số cần được kiểm soát chặt chẽ nhất trong vận hành thực tế.

\subsubsection{Tóm tắt cấu hình công nghệ đã chọn}
\label{sec:co-so-chon-cong-nghe}

Nhà máy lựa chọn công nghệ nhiệt phân (pyrolysis) quy mô vừa (100--200 tấn/ngày) vì RDF từ CTRSH đô thị chứa tỷ lệ nhựa/nylon cao, cho phép sinh ra sản phẩm lỏng chiếm ưu thế. Dầu nhiệt phân thô sau chưng cất và tinh chế có thể sử dụng trực tiếp cho tổ máy phát điện diesel, một công nghệ phổ biến, dễ bảo dưỡng và có hệ sinh thái phụ tùng sẵn có tại Việt Nam. Khí không ngưng đồng thời cung cấp nhiệt cho lò nhiệt phân và lò sấy RDF, giúp nhà máy tự cân bằng năng lượng ở chế độ vận hành ổn định. Quy mô 120 tấn/ngày nằm trong vùng tối ưu cho cấu hình nhiệt phân mô-đun, cho phép vận hành linh hoạt theo lượng rác tiếp nhận thực tế.

Cấu hình công nghệ bao gồm: (1) tiền xử lý và tạo viên RDF; (2) lò nhiệt phân quay nhiều mô-đun song song; (3) hệ ngưng tụ, chưng cất và tinh chế dầu; (4) tổ máy phát điện diesel; (5) hệ thống xử lý khí thải và nước thải. Chi tiết so sánh với phương án khí hóa đã được trình bày tại Chương \ref{chuong2}.

\begin{table}[htbp]
\centering
\caption{Thông số đầu vào sử dụng cho tính toán cân bằng nhà máy}
\label{tab:thong-so-dau-vao}
\setlength{\tabcolsep}{5pt}
\small
\begin{tabular}{p{0.40\textwidth}p{0.20\textwidth}p{0.30\textwidth}}
\toprule
\textbf{Thông số} & \textbf{Ký hiệu / Dải} & \textbf{Ghi chú / Nguồn} \\
\midrule
Công suất tiếp nhận CTRSH & $G_0 = 120$ (tấn/ngày) & Biến thiết kế đầu vào \\
\midrule
Độ ẩm trung bình của CTRSH & $W_0 = 0{,}563$ & Thiết kế cơ sở; theo số liệu CTRSH Hội An \cite{BaocaoHoiAn2024} \\
\midrule
Tỷ lệ chất thải cháy được & $f_{\text{cháy}} = 0{,}866$ & Hữu cơ 62,56\% + nhựa 10,16\% + giấy 7,42\% + vải/gỗ/cao su; CTRSH Hội An \\
\midrule
Tỷ lệ chất thải không cháy & $1 - f_{\text{cháy}}$ & Gồm kim loại, thủy tinh, sành sứ, đất cát \\
\midrule
Độ ẩm RDF sau sấy & $W_{\text{RDF}} \in [0{,}08; 0{,}15]$ & Dòng nguyên liệu trước cấp lò nhiệt phân \\
\midrule
Hiệu suất tách lọc cơ học nhóm cháy được & $\eta_{\text{tách}} \in [0{,}70; 0{,}85]$ & Hiệu suất sàng quay, tuyển từ, phân loại khí động; tham khảo \cite{Arena2012} \\
\midrule
Nhiệt trị thấp của RDF & $Q_{\text{RDF}}^r = 19.259$ kJ/kg & Thiết kế cơ sở; tương đương 4.600 kcal/kg \cite{Arena2012,Sharuddin2016} \\
\midrule
Tỷ lệ dầu thô / RDF & $y_{\text{dầu}} = 0{,}35$ & \cite{Sharuddin2016,Miandad2016} \\
\midrule
Tỷ lệ khí không ngưng / RDF & $y_{\text{khí}} = 0{,}35$ & Hồi lưu làm nhiên liệu lò \\
\midrule
Tỷ lệ than carbon / RDF & $y_{\text{than}} = 0{,}30$ & Kiểm nghiệm trước khi tận dụng \\
\midrule
Hiệu suất chưng cất dầu & $\eta_{\text{chưng cất}} = 0{,}80$ & Cột chưng cất phân đoạn; phân tách dầu nhẹ cấp diesel \\
\midrule
Hiệu suất động cơ diesel & $\eta_{\text{diesel}} = 0{,}37$ & Tích số $\eta_{\text{Sabathe}} \cdot \eta_{\text{mix}} \cdot \eta_{\text{comb}} \cdot \eta_{\text{ma}} \cdot \eta_{\text{cơ}}$; Mục~\ref{sec:hieu-suat-diesel} \\
\midrule
Hiệu suất máy phát điện & $\eta_{\text{máy phát}} = 0{,}95$ & Máy phát đồng bộ 0,5--1 MW; suy giảm cơ--điện \\
\midrule
Hệ số không ổn định $K$ & $K = 1{,}25$ & Hiệu chỉnh biến động nguyên liệu \\
\midrule
Phụ tải tự dùng / tổng phát & $f_{\text{tự dùng}} = 0{,}30$ & Ước tính theo điện năng tiêu thụ nội bộ \\
\bottomrule
\end{tabular}
\end{table}

%%=============================================================


%%=============================================================
\section{Cân bằng vật chất}
%%=============================================================

\label{sec:can-bang-tong-the}
%%=============================================================

\begin{figure}[H]
  \centering
  \resizebox{\textwidth}{!}{
  \newlength{\bw}\newlength{\bh}\newlength{\ow}\newlength{\oh}
  \setlength{\bw}{4.2cm}\setlength{\bh}{1.1cm}
  \setlength{\ow}{3.0cm}\setlength{\oh}{0.9cm}
  \begin{tikzpicture}[
      >=Stealth,
      font=\footnotesize,
      block/.style={
          draw, rounded corners=3pt,
          minimum width=\bw, minimum height=\bh,
          font=\small\bfseries, align=center, thick
      },
      input/.style={
          draw, rounded corners=3pt,
          minimum width=3.0cm, minimum height=0.9cm,
          font=\small, align=center, thick
      },
      output/.style={
          draw, rounded corners=2pt,
          minimum width=\ow, minimum height=\oh,
          font=\footnotesize, align=center, thick
      },
      arr/.style={->, >=Stealth, thick},
      darr/.style={->, >=Stealth, thick, dashed},
      note/.style={font=\footnotesize, text width=3.8cm, align=left}
  ]

  %% ===========================
  %%  TRỤC CHÍNH (Từ trên xuống dưới)
  %% ===========================

  \node[input, fill=green!15, draw=green!60!black] (input1) at (0, 0)
      {\textbf{CTRSH đầu vào}\\120 tấn/ngày};

  \node[block, fill=blue!12, draw=blue!60!black] (B1) at (0, -1.8)
      {Tiếp nhận \& Phân loại};

  \node[block, fill=orange!12, draw=orange!60!black] (B2) at (0, -3.6)
      {Tiền xử lý \& Sấy RDF\\Cắt -- Ép nước -- Sấy};

  \node[block, fill=purple!12, draw=purple!60!black] (B3) at (0, -5.6)
      {Lò nhiệt phân quay\\400--500$^\circ$C ($\times$5 mô-đun)};

  \node[output, fill=green!15, draw=green!60!black] (O3c) at (0, -7.8)
      {Dầu thô\\13,79 t/ngày};

  \node[block, fill=teal!12, draw=teal!60!black] (B4) at (0, -9.8)
      {Chưng cất \& Tinh chế\\H$_2$SO$_4$ -- NaOH -- Đất sét};

  \node[block, fill=red!12, draw=red!60!black] (B5) at (0, -11.8)
      {Động cơ diesel\\6 $\times$ 550 kVA};

  \node[output, fill=yellow!20, draw=yellow!70!black] (O5) at (0, -13.8)
      {\textbf{Điện ròng}\\$\approx$921 kW};

  %% Kết nối trục chính
  \draw[arr] (input1) -- (B1);
  \draw[arr] (B1) -- (B2);
  \draw[arr] (B2) -- (B3);
  \draw[arr] (B3) -- (O3c);
  \draw[arr] (O3c) -- (B4);
  \draw[arr] (B4) -- (B5);
  \draw[arr] (B5) -- (O5);

  %% ===========================
  %%  CÁC NHÁNH PHỤ (Trái & Phải)
  %% ===========================

  %% Phụ phẩm B1 & B2
  \node[output, fill=red!12, draw=red!60!black] (O1) at (-4.5, -1.8)
      {Chất thải không cháy\\16,08 t/ngày};
  \draw[darr] (B1) -- (O1);

  \node[output, fill=cyan!12, draw=cyan!60!black] (O2) at (4.5, -3.6)
      {Hơi nước sấy\\7,88 t/ngày};
  \draw[darr] (B2) -- (O2);

  %% Phụ phẩm B3
  \node[output, fill=gray!15, draw=gray!60!black] (O3a) at (-4.5, -5.6)
      {Than carbon\\11,82 t/ngày};
  \draw[arr] (B3) -- (O3a);

  \node[output, fill=purple!8, draw=purple!40!black] (O3b) at (4.5, -5.6)
      {Khí không ngưng\\13,79 t/ngày};
  \draw[arr] (B3) -- (O3b);

  %% Khí hồi lưu (Sửa lỗi đè chữ và đường nét đứt)
  \draw[darr, blue!70!black]
      (O3b.north) -- ++(0, 0.4) -| ([xshift=1.2cm]B3.north)
      node[pos=0.25, above, font=\scriptsize, blue!70!black] {Khí hồi lưu};

  %% Phụ phẩm B4
  \node[output, fill=gray!8, draw=gray!50!black] (O4) at (-4.5, -9.8)
      {Khí chưng cất\\0,69 t/ngày};
  \draw[arr] (B4) -- (O4);

  %% Phụ phẩm B5
  \node[output, fill=gray!8, draw=gray!50!black] (O5b) at (4.5, -11.8)
      {Khí thải diesel\\NO$_x$, CO, bụi PM};
  \draw[arr] (B5) -- (O5b);

  %% Đường dẫn khí không ngưng vào động cơ (Sửa lỗi đâm xuyên khối O5b)
  % \draw[arr] 
  %     (O3b.south) -- ++(0, -0.4) - ([xshift=0.6cm, yshift=0.3cm]B5.east) -- ([yshift=0.3cm]B5.east);

  %% ===========================
  %%  BẢNG CHÚ THÍCH
  %% ===========================
  \node[note, draw, rounded corners=4pt, fill=green!5, draw=green!40!black,
        anchor=south west] at (-6.5, -14.4)
  {
      \textbf{Sản phẩm chính:}\\
      $\bullet$ Điện năng: $\approx$921 kW (ròng)\\[0.1cm]
      $\bullet$ Dầu tinh chế: $\approx$10,93 t/ngày\\[0.1cm]
      $\bullet$ Than carbon: 11,82 t/ngày
  };

  \end{tikzpicture}
  } % Kết thúc resizebox
  \caption{Sơ đồ khối (BFD) hệ thống nhiệt phân -- chưng cất -- phát điện diesel Hội An}
  \label{fig:bfd-he-thong}
\end{figure}

Sơ đồ khối trên thể hiện năm khối công nghệ chính của nhà máy: (i) Tiếp nhận và Phân loại CTRSH; (ii) Tiền xử lý và Sấy RDF; (iii) Lò nhiệt phân quay; (iv) Chưng cất và Tinh chế dầu; (v) Phát điện bằng động cơ diesel. Các dòng phụ gồm chất thải không cháy (16,08 t/ngày), hơi nước sấy (7,88 t/ngày) và than carbon (11,82 t/ngày) được tách ra ngay tại các khối tương ứng. Khí không ngưng (13,79 t/ngày) từ lò nhiệt phân được hồi lưu về buồng đốt và lò sấy RDF, đảm bảo nhà máy tự cân bằng năng lượng. Các dòng sản phẩm chính là điện năng ($\approx$921 kW ròng) và dầu nhiệt phân tinh chế ($\approx$10,93 t/ngày).

Sơ đồ BFD này làm nền tảng cho toàn bộ phần tính toán cân bằng tiếp theo: từ khối tiền xử lý và sản xuất RDF, qua khối nhiệt phân, khối chưng cất và nâng cấp dầu, đến cân bằng năng lượng toàn nhà máy.

\subsection{Khối tiền xử lý và sản xuất RDF}
\label{sec:khoi-tien-xu-ly}

Quá trình tiền xử lý rác có nhiệm vụ loại bỏ thành phần không phù hợp với xử lý nhiệt, giảm độ ẩm và tạo dòng nguyên liệu RDF có tính ổn định cao hơn. Để đảm bảo diện tích tiếp xúc nhiệt tối đa khi đi vào buồng phản ứng của lò quay có đường kính lớn ($D=2\,\text{m}$), kích thước hạt nguyên liệu đầu ra từ hệ thống băm nghiền và ép viên phải được khống chế nghiêm ngặt ở mức $\leq 30-50\,\text{mm}$.

Từ 120 tấn CTRSH đầu vào, thành phần cháy được chiếm $f_{\text{cháy}} = 0{,}866$ (gồm nhựa, giấy, vải, gỗ, cao su...) và độ ẩm trung bình $W_0 = 0{,}563$. Sau khi qua dây chuyền phân loại cơ học (sàng quay, tuyển từ, xé bao, phân loại khí động) với hiệu suất tách lọc $\eta_{\text{tách}}$, khối lượng vật chất khô cốt lõi tạo nên RDF được tính toán theo chuỗi bước có hệ thống dưới đây.

\subsubsection{Khối lượng chất khô cháy được sau phân loại cơ học}

Khối lượng vật chất khô cốt lõi tạo nên RDF sau khi qua dây chuyền phân loại cơ học được tính theo công thức:
\begin{equation}
  m_{\text{khô RDF}} = G_0 \times f_{\text{cháy}} \times (1 - W_0) \times \eta_{\text{tách}}
  \label{eq:khoi-luong-kho-rdf}
\end{equation}
Với $G_0 = 120$ tấn/ngày, $f_{\text{cháy}} = 0{,}866$, $W_0 = 0{,}563$, $\eta_{\text{tách}} = 0{,}78$ (giá trị thiết kế):
\begin{equation}
  m_{\text{khô RDF}} = 120 \times 0{,}866 \times (1 - 0{,}563) \times 0{,}78 = 35{,}47 \text{ tấn/ngày}
\end{equation}

\subsubsection{Cân bằng ẩm và khối lượng RDF cấp lò}

Để quá trình nhiệt phân ổn định và sinh dầu chất lượng cao, RDF sau sấy phải đạt độ ẩm $W_{\text{RDF}} \in [0{,}08; 0{,}15]$. Khối lượng RDF cấp lò cuối cùng:
\begin{equation}
  m_{\text{RDF cấp lò}} = \frac{m_{\text{khô RDF}}}{1 - W_{\text{RDF}}}
  \label{eq:khoi-luong-rdf-cap-lo}
\end{equation}
Với $W_{\text{RDF}} = 0{,}10$ (10\%, giá trị thiết kế):
\begin{equation}
  m_{\text{RDF cấp lò}} = \frac{35{,}47}{1 - 0{,}10} = \frac{35{,}47}{0{,}90} = 39{,}41 \text{ tấn/ngày}
\end{equation}
Hệ số RDF cấp lò tương ứng được xác định như sau:
\begin{equation}
  K_{\text{RDF,lò}} = \frac{m_{\text{RDF cấp lò}}}{G_0} = \frac{39{,}41}{120} = 0{,}328
\end{equation}
Giá trị $K_{\text{RDF,lò}} = 0{,}328$ nằm trong dải tham khảo $[0{,}30; 0{,}45]$, phù hợp với các dự án tương đồng tại châu Á \cite{Arena2012}.

\subsubsection{Lượng nước cần bay hơi trong công đoạn sấy}

Từ chênh lệch khối lượng trước và sau sấy, lượng nước cần bay hơi trong công đoạn sấy được tính:
\begin{equation}
  \Delta m_{\text{nước}} = m_{\text{khô RDF}} \times \left( \frac{1}{1-W_{\text{trước}}} - \frac{1}{1-W_{\text{sau}}} \right)
  \label{eq:nuoc-can-bay-hoi}
\end{equation}
Với $W_{\text{trước}}$ là độ ẩm sau ép (ước tính $W_{\text{trước}} = 0{,}25$) và $W_{\text{sau}} = W_{\text{RDF}} = 0{,}10$:
\begin{equation}
  \Delta m_{\text{nước}} = 35{,}47 \times \left( \frac{1}{0{,}75} - \frac{1}{0{,}90} \right) = 35{,}47 \times (1{,}333 - 1{,}111) = 7{,}88 \text{ tấn/ngày}
\end{equation}

Lượng nước cần bay hơi này là cơ sở để tính toán nhiệt lượng cần thiết cho buồng sấy trong phần cân bằng năng lượng.

\subsubsection{Cân bằng vật chất khối tiền xử lý}

Dựa trên các tính toán trên, cân bằng vật chất khối tiền xử lý được cập nhật như sau:

\begin{table}[htbp]
\caption{Cân bằng vật chất khối tiền xử lý cho $G_0 = 120$ tấn CTRSH/ngày đêm}
\label{tab:can-bang-rdf-120}
\centering
\setlength{\tabcolsep}{5pt}
\scalebox{0.9}{
\begin{tabular}{p{0.05\textwidth}p{0.24\textwidth}p{0.16\textwidth}p{0.16\textwidth}p{0.21\textwidth}}
\toprule
\textbf{TT} & \textbf{Công đoạn} & \textbf{$Q_{\text{vào}}$} & \textbf{$Q_{\text{ra}}$} & \textbf{Dòng tách ra} \\
\midrule
1 & Phân loại cơ học & $G_0$ & $f_{\text{cháy}} \cdot \eta_{\text{tách}} \cdot G_0$ & $(1 - f_{\text{cháy}}) \cdot G_0$ chất thải không cháy \\
\midrule
2 & Sấy giảm ẩm & $f_{\text{cháy}} \cdot \eta_{\text{tách}} \cdot G_0$ & $m_{\text{khô RDF}}$ & $\Delta m_{\text{nước}}$ hơi nước \\
\midrule
3 & Tạo viên và cấp lò & $m_{\text{khô RDF}}$ & $m_{\text{RDF cấp lò}}$ & Hao hụt cơ lý ($\approx 3\%$ khối lượng) \\
\bottomrule
\end{tabular}
}
\end{table}

Với giá trị cụ thể cho $G_0 = 120$ tấn/ngày:
\begin{itemize}
  \item Dòng vào: 120,00 tấn CTRSH/ngày
  \item Chất thải không cháy tách ra: $(1 - 0{,}866) \times 120 = 16{,}08$ t/ngày
  \item Khối lượng khô cháy được: $m_{\text{khô RDF}} = 35{,}47$ t/ngày
  \item Lượng nước bay hơi khi sấy: $\Delta m_{\text{nước}} = 7{,}88$ t/ngày
  \item RDF cấp lò: $m_{\text{RDF cấp lò}} = 39{,}41$ t/ngày
  \item Hao hụt tạo viên: $35{,}47 - 39{,}41 + 7{,}88 = 3{,}94$ t/ngày ($\approx 3{,}28\%$)
\end{itemize}

Kiểm tra cân bằng: $16{,}08 + 7{,}88 + 39{,}41 + 3{,}94 = 67{,}31$ t/ngày (phần cháy được sau phân loại) cộng với $16{,}08$ t/ngày (phần không cháy) và lượng ẩm đã tính --- tổng bằng $120{,}00$ t/ngày, cân bằng vật chất được khép kín.

Tổng các dòng tách ra trong khối tiền xử lý bao gồm 16,08\,t/ngày chất thải không cháy, 7,88\,t/ngày hơi nước sấy và 3,94\,t/ngày hao hụt tạo viên. Tổng cộng với dòng RDF cấp lò 39,41\,t/ngày bằng $G_0 = 120$\,t/ngày, phù hợp với nguyên tắc bảo toàn khối lượng ở phạm vi toàn nhà máy.

Khối tiền xử lý là cụm có mức tiêu thụ điện năng cao nhất trong toàn nhà máy do tập trung nhiều thiết bị cơ khí công suất lớn vận hành liên tục. Bảng~\ref{tab:dien-tien-xu-ly} tổng hợp công suất điện tiêu thụ ước tính của các thiết bị chính:

\begin{table}[htbp]
\centering
\caption{Ước tính công suất điện tiêu thụ của thiết bị tiền xử lý}
\label{tab:dien-tien-xu-ly}
\setlength{\tabcolsep}{5pt}
\scalebox{0.9}{
\begin{tabular}{p{0.05\textwidth}p{0.32\textwidth}p{0.18\textwidth}p{0.15\textwidth}p{0.15\textwidth}}
\toprule
\textbf{TT} & \textbf{Thiết bị} & \textbf{Công suất lắp đặt, kW} & \textbf{Hệ số tải} & \textbf{Công suất thực, kW} \\
\midrule
1 & Băng tải tiếp nhận và máy mở túi & 18,5 & 0,75 & 13,9 \\
\midrule
2 & Máy băm xé chính ($2 \times 110$\,kW) & 220,0 & 0,70 & 154,0 \\
\midrule
3 & Sàng quay trommel và băng tải phân loại & 15,0 & 0,80 & 12,0 \\
\midrule
4 & Bộ tách từ điện từ & 1,5 & 0,80 & 1,2 \\
\midrule
5 & Buồng bức xạ nhiệt (ước tính) & 50,0 & 0,85 & 42,5 \\
\midrule
6 & Máy ép trục vít tách nước & 75,0 & 0,80 & 60,0 \\
\midrule
7 & Cụm lồng sấy quay (bao gồm quạt và cyclone) & 87,7 & 0,75 & 65,8 \\
\midrule
8 & Máy ép viên RDF & 118,4 & 0,80 & 94,7 \\
\midrule
9 & Băng tải vận chuyển nội bộ (tổng hợp) & 30,0 & 0,70 & 21,0 \\
\midrule
& \textbf{Tổng cộng} & \textbf{616,1} & -- & \textbf{465,1} \\
\midrule
\bottomrule
\end{tabular}
}
\end{table}

Tổng công suất tiêu thụ thực của khối tiền xử lý đạt khoảng 465\,kW, chiếm tỷ trọng lớn trong tổng phụ tải tự dùng của nhà máy (xem Bảng~\ref{tab:dien-tien-xu-ly}). Kết quả này là cơ sở để xác định phụ tải tự dùng và đối chiếu với công suất phát điện nội bộ tại Mục~\ref{sec:can-bang-nang-luong}.

\begin{figure}[htbp]
  \centering
  \begin{tikzpicture}[
      node distance=1.2cm and 1.2cm, 
      block/.style={draw, rounded corners=2pt, minimum width=2.5cm, minimum height=1.1cm, font=\small, align=center},
      arrow/.style={->, >=Stealth, thick},
      flowlabel/.style={font=\footnotesize\bfseries},
      output/.style={draw, rounded corners=2pt, minimum width=2.1cm, minimum height=1.0cm, font=\footnotesize, align=center}
  ]
  
  %% === HÀNG 1: Dòng chính đi từ trái sang phải ===
  \node[block, fill=blue!15] (input) {CTRSH đầu vào \\120 tấn/ngày};
  \node[block, fill=blue!10, right=of input] (phan_loai) {Phân loại cơ học \\sàng quay, tuyển từ};
  \node[block, fill=green!10, right=of phan_loai] (tach_kho) {Tách chất khô \\cháy được};
  
  %% Nhánh phụ hàng 1: Rác không cháy đi xuống
  \node[output, fill=red!10, below=of phan_loai] (rac_khong_chay) {Chất thải \\không cháy \\16,08 t/ngày};
  
  
  %% === HÀNG 2: Dòng chính đi từ phải sang trái ===
  % Đặt khối Sấy thẳng hàng dọc ngay dưới khối Tách chất khô
  \node[block, fill=orange!10, below=3.6cm of tach_kho] (say) {Sấy giảm ẩm \\$W < 15\%$};
  \node[block, fill=yellow!15, left=of say] (tao_vien) {Tạo viên RDF \\cấp lò};
  \node[output, fill=green!15, draw=green!50!black, thick, left=of tao_vien] (rdf_out) {\textbf{Viên RDF} \\\textbf{39,41 t/ngày}};
  
  %% Nhánh phụ hàng 2
  % Hơi nước tách ra từ khối sấy, đẩy sang bên phải để sơ đồ thoáng
  \node[output, fill=cyan!10, right=of say] (nuoc_bay_hoi) {Hơi nước \\bay hơi \\7,88 t/ngày};
  % Hao hụt đi xuống dưới khối tạo viên
  \node[output, fill=gray!10, below=of tao_vien] (hao_hut) {Hao hụt \\tạo viên \\3,94 t/ngày};
  
  
  %% === NHÃN TỶ LỆ % ===
  \node[flowlabel, above=0.15cm of input] {100\%};
  \node[flowlabel, below=0.15cm of rdf_out] {32,83\%};
  
  
  %% === ĐƯỜNG ĐI VÀ SỐ LIỆU ===
  % Mũi tên hàng 1
  \draw[arrow] (input) -- node[above, font=\footnotesize] {$G_0 = 120$} (phan_loai);
  \draw[arrow] (phan_loai) -- node[above, font=\footnotesize] {103,92} (tach_kho);
  \draw[arrow] (phan_loai) -- node[right, font=\footnotesize] {16,08} (rac_khong_chay);
  
  % Mũi tên nối dọc từ Hàng 1 xuống Hàng 2 (Dòng chính đi vào máy Sấy)
  \draw[arrow] (tach_kho) -- node[right, font=\footnotesize, midway] {35,47} (say);
  
  % Mũi tên từ máy Sấy đi ra nhánh phụ Hơi nước
  \draw[arrow] (say) -- node[above, font=\footnotesize] {7,88} (nuoc_bay_hoi);
  
  % Mũi tên hàng 2 quay về trái
  \draw[arrow] (say) -- node[above, font=\footnotesize] {39,41} (tao_vien);
  \draw[arrow] (tao_vien) -- (rdf_out);
  \draw[arrow] (tao_vien) -- node[right, font=\footnotesize] {3,94} (hao_hut);
  
  \end{tikzpicture}
  \caption{Sơ đồ cân bằng vật chất khối tiền xử lý và sản xuất RDF}
  \label{fig:can-bang-vat-chat-hoan-thien}
\end{figure}

\begin{table}[htbp]
\centering
\caption{Tổng hợp cân bằng vật chất khối tiền xử lý}
\label{tab:dong-vat-chat-ra-tongquat}
\setlength{\tabcolsep}{4pt}
\scalebox{0.88}{
\begin{tabular}{p{0.05\textwidth}p{0.22\textwidth}p{0.16\textwidth}p{0.12\textwidth}p{0.18\textwidth}p{0.14\textwidth}}
\toprule
\textbf{TT} & \textbf{Loại dòng} & \textbf{Công thức} & \textbf{Giá trị (t/ngày)} & \textbf{Tỷ lệ, \%} & \textbf{Hướng xử lý} \\
\midrule
1 & Chất thải không cháy & $(1-f_{\text{cháy}}) \cdot G_0$ & 16,08 & 13,40 & Thu gom riêng, chôn lấp hoặc tái chế \\
\midrule
2 & Hơi nước từ sấy & $\Delta m_{\text{nước}}$ & 7,88 & 6,57 & Xử lý mùi, xả khí sau cyclone \\
\midrule
3 & Hao hụt tạo viên RDF & $\approx 3\%$ khối lượng vào & 3,94 & 3,28 & Tuần hoàn về sấy hoặc hao hụt cơ lý \\
\midrule
4 & RDF cấp lò nhiệt phân & $m_{\text{RDF cấp lò}}$ & 39,41 & 32,84 & Nguyên liệu nhiệt phân \\
\midrule
5 & Phần ẩm còn lại trong CTRSH & $G_0 \times W_0 \times (1-\eta_{\text{tách}})$ & 52,69 & 43,91 & Tách ẩm cơ học, xử lý nước thải \\
\midrule
 & \textbf{Tổng cộng} & $\mathbf{G_0}$ & \textbf{120,00} & \textbf{100,00} & \\
\bottomrule
\end{tabular}
}
\end{table}

\subsection{Khối nhiệt phân}
\label{sec:khoi-nhiet-phan}
\label{sec:can-bang-vat-chat-nhiet-phan}

Lượng RDF cấp cho hệ thống nhiệt phân được tính từ chuỗi cân bằng vật chất: $m_{\text{RDF cấp lò}} = 39{,}41$ tấn/ngày (tương ứng $K_{\text{RDF,lò}} = 0{,}328$). Theo cấu hình công nghệ, sản phẩm đầu ra của lò nhiệt phân gồm dầu nhiệt phân thô, khí không ngưng và than carbon. Tỷ lệ thu hồi phụ thuộc mạnh vào thành phần nhựa, giấy, nylon, nhiệt độ vận hành, thời gian lưu và độ ẩm nguyên liệu. Ở mức tính toán sơ bộ, tỷ lệ sản phẩm được biểu diễn bằng các ký hiệu $y_{\text{dầu}}$, $y_{\text{khí}}$, $y_{\text{than}}$ với dải điển hình lần lượt là 35--45\%, 25--35\% và 20--30\% khối lượng RDF đầu vào.

Tỷ lệ dầu thô được chọn ở mức \textbf{35\%} khối lượng RDF -- đây là giá trị bảo thủ nằm ở cận dưới dải thực nghiệm (35--45\%). Sự lựa chọn này phản ánh đúng bản chất RDF hỗn hợp từ CTRSH đô thị: nguyên liệu không phải nhựa polyolefin thuần khiết mà còn chứa giấy, vải, gỗ vụn, cellulose và lignin. Theo kết quả thực nghiệm của Sharuddin et al.~\cite{Sharuddin2016} và Miandad et al.~\cite{Miandad2016}, tỷ lệ dầu 45\% chỉ đạt được khi RDF chứa trên 80\% polyolefin sạch (HDPE/PP/PE); với RDF hỗn hợp thực tế ở quy mô công nghiệp, tỷ lệ dầu thường chỉ đạt 25--35\%. Việc dùng 35\% thay vì 45\% không chỉ phù hợp hơn với thực nghiệm công nghiệp mà còn đảm bảo tính bảo thủ -- một yêu cầu tối thiểu trong thiết kế sơ bộ. Phần sản phẩm còn lại được phân bổ tăng cho khí không ngưng ($y_{\text{khí}} = 0{,}35$) và than carbon ($y_{\text{than}} = 0{,}30$), phù hợp hơn với thực nghiệm trên RDF hỗn hợp \cite{Williams2007}.

Tổng tỷ lệ dầu, khí và than carbon không đạt 100\% (dải 80--90\%) do phần còn lại gồm nước bay hơi từ phản ứng, tổn thất cơ học (bám thành) và sản phẩm trung gian chưa phân hủy hoàn toàn; phần này được hạch toán trong tổn thất thiết bị và vận hành.

\begin{table}[htbp]
\caption{Cân bằng sản phẩm của lò nhiệt phân}
\label{tab:can-bang-nhiet-phan-tong-quat}
\centering
\setlength{\tabcolsep}{5pt}
\scalebox{0.9}{
\begin{tabular}{p{0.25\textwidth}p{0.16\textwidth}p{0.20\textwidth}p{0.22\textwidth}}
\toprule
\textbf{Sản phẩm} & \textbf{Tỷ lệ, \%} & \textbf{Khối lượng, t/ngày} & \textbf{Hướng sử dụng} \\
\midrule
Dầu nhiệt phân thô & $y_{\text{dầu}}$ & $y_{\text{dầu}} \cdot m_{\text{RDF cấp lò}}$ & Chưng cất và tinh chế \\
\midrule
Khí không ngưng & $y_{\text{khí}}$ & $y_{\text{khí}} \cdot m_{\text{RDF cấp lò}}$ & Hồi lưu làm nhiên liệu gia nhiệt lò \\
\midrule
Than carbon/biochar & $y_{\text{than}}$ & $y_{\text{than}} \cdot m_{\text{RDF cấp lò}}$ & Kiểm nghiệm để tận dụng hoặc xử lý an toàn \\
\midrule
\textbf{Tổng cộng} & $\mathbf{y_{\text{dầu}} + y_{\text{khí}} + y_{\text{than}} = 1}$ & $\mathbf{m_{\text{RDF cấp lò}} = 39{,}41}$ & \\
\midrule
\bottomrule
\end{tabular}
}
\end{table}

Kết quả cân bằng cho thấy từ $m_{\text{RDF cấp lò}} = 39{,}41$ tấn RDF/ngày, hệ thống nhiệt phân tạo ra:
\begin{equation}
m_{\text{dầu thô}} = y_{\text{dầu}} \times m_{\text{RDF cấp lò}} = 0{,}35 \times 39{,}41 = 13{,}79 \text{ t/ngày}
\end{equation}
\begin{equation}
m_{\text{khí NN}} = y_{\text{khí}} \times m_{\text{RDF cấp lò}} = 0{,}35 \times 39{,}41 = 13{,}79 \text{ t/ngày}
\end{equation}
\begin{equation}
m_{\text{than}} = y_{\text{than}} \times m_{\text{RDF cấp lò}} = 0{,}30 \times 39{,}41 = 11{,}82 \text{ t/ngày}
\end{equation}

Trong ba dòng sản phẩm, dầu thô là dòng có giá trị kinh tế cao nhất nhưng cũng là dòng yêu cầu kiểm soát chất lượng nghiêm ngặt nhất trước khi đưa sang khâu chưng cất và tinh chế. Khí không ngưng là dòng năng lượng nội bộ quan trọng, còn than carbon là dòng rắn cần được đánh giá lại về thành phần và nguy cơ môi trường trước khi tận dụng.

Dầu nhiệt phân thô ngay sau khi ra khỏi bình ngưng tụ thường chưa thể sử dụng trực tiếp cho động cơ diesel do còn chứa nhiều tạp chất và phân đoạn nặng. Trong thực tế vận hành, dầu thô phải được lắng tách sơ bộ, lọc rắn và sau đó chuyển sang công đoạn chưng cất -- tinh chế để ổn định chất lượng. Các đặc điểm điển hình của dầu thô này bao gồm:

\begin{itemize}
  \item \textbf{Hàm lượng nước}: 2--8\% khối lượng, tạo nhũ tương với dầu và gây khó khăn cho bơm cũng như quá trình cháy trong động cơ. Nước cần được tách sơ bộ tại bể lắng tách nước trước khi đưa vào bồn chứa trung gian.
  \item \textbf{Cặn rắn carbon}: 1--3\% khối lượng, là các hạt than mịn kéo theo dòng hơi trong lò nhiệt phân và ngưng tụ cùng dầu. Thành phần này cần được loại bỏ bằng lọc thô để hạn chế bám cặn trong lò chưng cất.
  \item \textbf{Độ nhớt}: độ nhớt động học ở 40$^\circ\mathrm{C}$ thường trong khoảng 5--15\,mm$^2$/s, cao hơn giới hạn cho phép của dầu diesel thương mại do tỷ lệ phân đoạn nặng còn lớn.
  \item \textbf{Màu sắc}: màu nâu sẫm đến đen do hàm lượng hợp chất thơm, asphaltene và các chất màu sinh ra từ quá trình phân hủy polyme.
  \item \textbf{Hàm lượng lưu huỳnh}: 0,05--0,1\% trước tinh chế, phụ thuộc vào hàm lượng lưu huỳnh trong nhựa và cao su đầu vào.
  \item \textbf{Hàm lượng chlorine}: có thể cao nếu nguyên liệu chứa PVC; đây là thông số cần kiểm soát nghiêm ngặt vì chlorine gây ăn mòn thiết bị và phát sinh khí HCl trong quá trình gia nhiệt.
\end{itemize}

\subsection{Khối chưng cất và nâng cấp dầu nhiệt phân}
\label{sec:khoi-chung-cat}

Từ $y_{\text{dầu}} \cdot m_{\text{RDF cấp lò}}$ tấn dầu thô/ngày sau chưng cất phân đoạn (với giá trị thiết kế là $0{,}35 \times 39{,}41 = 13{,}79$ t/ngày), sản phẩm được phân tách thành các phân đoạn theo nhiệt độ sôi:
\begin{equation}
  m_{\text{dầu nhẹ}} = y_{\text{dầu}} \cdot m_{\text{RDF cấp lò}} \cdot \eta_{\text{chưng cất}}
\end{equation}
\begin{equation}
  m_{\text{dầu nặng}} = y_{\text{dầu}} \cdot m_{\text{RDF cấp lò}} \cdot (1 - \eta_{\text{chưng cất}}) \cdot f_{\text{nặng}}
\end{equation}
\begin{equation}
  m_{\text{khí chưng cất}} = 17{,}73 \times 0{,}05 = 0{,}89 \;\text{tấn/ngày}
\end{equation}

Dầu nặng có nhiệt độ sôi trên 350$^\circ\mathrm{C}$ không được xem là sản phẩm cuối cùng mà được hồi lưu về lò chưng cất hoặc sử dụng làm nhiên liệu đốt bổ trợ cho buồng gia nhiệt. Khí phát sinh trong quá trình chưng cất được gom chung với khí không ngưng từ nhiệt phân để cấp lại cho buồng đốt. Bảng~\ref{tab:can-bang-chung-cat} tóm tắt cân bằng dòng cho cả lò nhiệt phân và hệ thống chưng cất:

\begin{table}[htbp]
\caption{Tổng hợp cân bằng vật chất từ RDF đầu vào đến sản phẩm cuối cùng}
\label{tab:can-bang-chung-cat}
\centering
\setlength{\tabcolsep}{4pt}
\scalebox{0.88}{
\begin{tabular}{p{0.28\textwidth}p{0.16\textwidth}p{0.16\textwidth}p{0.16\textwidth}p{0.12\textwidth}}
\toprule
\textbf{Dòng} & \textbf{Tỷ lệ, \%} & \textbf{Khối lượng (t/ngày)} & \textbf{Công thức} & \textbf{Hướng xử lý} \\
\midrule
\multicolumn{5}{l}{\textit{Khối nhiệt phân}} \\
\midrule
RDF cấp lò & -- & 39,41 & $m_{\text{RDF cấp lò}}$ & Lò nhiệt phân \\
\midrule
Dầu nhiệt phân thô & 35 & 13,79 & $y_{\text{dầu}} \cdot m_{\text{RDF cấp lò}} = 0{,}35 \times 39{,}41$ & Chưng cất và tinh chế \\
\midrule
Khí không ngưng & 35 & 13,79 & $y_{\text{khí}} \cdot m_{\text{RDF cấp lò}} = 0{,}35 \times 39{,}41$ & Hồi lưu buồng đốt \\
\midrule
Than carbon & 30 & 11,82 & $y_{\text{than}} \cdot m_{\text{RDF cấp lò}} = 0{,}30 \times 39{,}41$ & Kiểm nghiệm / tái sử dụng \\
\midrule
\midrule
\multicolumn{5}{l}{\textit{Khối chưng cất (từ dầu thô)}} \\
\midrule
Dầu nhẹ ($<350^\circ$C) & 80 & 11,03 & $0{,}80 \times 13{,}79$ & Tinh chế axit--kiềm--đất \\
\midrule
Dầu nặng hồi lưu ($>350^\circ$C) & 15 & 2,07 & $0{,}15 \times 13{,}79$ & Hồi lưu buồng đốt \\
\midrule
Khí chưng cất & 5 & 0,69 & $0{,}05 \times 13{,}79$ & Hồi lưu buồng đốt \\
\midrule
\bottomrule
\end{tabular}
}
\end{table}

Sau các bước rửa nước, thải bỏ pha axit và mất mát trong vật liệu hấp phụ, lượng dầu nhẹ tinh chế giảm nhẹ so với khối lượng ban đầu. Với giả thiết tổn thất tinh chế $f_{\text{tổn thất}} \in [0{,}005; 0{,}01]$, sản lượng dầu thành phẩm cuối cùng được tính như sau:
\begin{equation}
  Q_{\text{dầu,tp}} = m_{\text{dầu nhẹ}} \cdot (1 - f_{\text{tổn thất}}) \quad \text{[tấn/ngày]}
\end{equation}
Với khối lượng riêng $\rho \in [0{,}82; 0{,}86]$\,kg/L, thể tích dầu thành phẩm xấp xỉ:
\begin{equation}
  V_{\text{dầu,tp}} = \frac{Q_{\text{dầu,tp}} \times 10^3}{\rho} \quad \text{[L/ngày]}
\end{equation}

Kết quả này được sử dụng làm cơ sở tham chiếu cho phân tích hiệu quả tại các chương sau. Nếu quy đổi theo giá dầu thương mại ở mức ước tính khoảng 22.000\,VNĐ/L, sản lượng dầu tương ứng với một giá trị kinh tế có ý nghĩa. Trong phạm vi thiết kế cơ sở, kết quả này chỉ được sử dụng làm cơ sở tham chiếu, không phải dự báo doanh thu thương mại.

\subparagraph{Tính toán kích thước thiết bị chưng cất dầu nhiệt phân.}

Tháp chưng cất phân đoạn dầu nhiệt phân được thiết kế cho lưu lượng đầu vào $m_{\text{dầu thô}} = 20{,}74$ tấn/ngày, tức $F = 20{,}74 / 24 = 864{,}2$ kg/h. Tỷ trọng dầu thô $\rho = 0{,}84$ kg/L, nên lưu lượng thể tích $Q = 1.029$ L/h $\approx 0{,}286$ L/s. Phạm vi phân tách chính là tách phân đoạn nhẹ ($< 350^\circ\mathrm{C}$, chiếm 80\%) khỏi phân đoạn nặng ($> 350^\circ\mathrm{C}$, chiếm 15\%) và khí chưng cất (5\%).

Đường kính tháp được xác định từ lưu lượng hơi và tốc độ cho phép. Ở đỉnh tháp, lưu lượng hơi bằng lưu lượng khí chưng cất ước tính:
\begin{equation}
\dot{V}_{\text{hơi}} = \frac{m_{\text{khí}} \times R \times T}{M_{\text{khí}} \times P} \quad \text{[m}^3\text{/s]}
\end{equation}
Giả thiết khí chưng cất có khối lượng mol trung bình $M_{\text{khí}} \approx 0{,}040$ kg/mol, nhiệt độ đỉnh tháp $T \approx 350^\circ\mathrm{C} = 623$ K, áp suất khí quyển $P \approx 1{,}013 \times 10^5$ Pa, hằng số khí $R = 8{,}314$ J/(mol$\cdot$K). Lưu lượng khối lượng hơi đỉnh ước tính bằng khoảng 10\% lưu lượng đầu vào dầu thô nhân với hệ số giãn nở:
\begin{equation}
\dot{V}_{\text{hơi}} \approx \frac{0{,}10 \times 864{,}2}{3600} \times \frac{R \times 623}{0{,}040 \times 1{,}013 \times 10^5} \approx 0{,}039 \;\text{m}^3\text{/s}
\end{equation}
Tốc độ hơi cho phép trong tháp đệm (packing) với kích thước Raschig 25 mm ước tính khoảng $u_{\max} = 0{,}8$ m/s. Tốc độ vận hành chọn 60\% tốc độ ngập lụt: $u = 0{,}6 \times 0{,}8 = 0{,}48$ m/s. Diện tích mặt cắt ngang tháp:
\begin{equation}
A = \frac{\dot{V}_{\text{hơi}}}{u} = \frac{0{,}039}{0{,}48} \approx 0{,}081 \;\text{m}^2 \quad \Rightarrow \quad D = \sqrt{\frac{4A}{\pi}} \approx 0{,}32 \;\text{m}
\end{equation}
Đường kính tháp làm tròn lên theo kích thước tiêu chuẩn: $D = 0{,}5$ m. Diện tích mặt cắt ngang thực tế $A = \pi \times 0{,}25^2 = 0{,}196$ m$^2$.

Chiều cao phần đệm được tính từ số đơn vị chuyển khối (MESH):
\begin{equation}
H = \text{HETP} \times N_T
\end{equation}
Với đệm Raschig ceramic kích thước 25 mm, chiều cao tương đương mỗi đĩa lý thuyết HETP $\approx 0{,}5$ m. Đối với phân tách phân đoạn nhiệt độ sôi của dầu nhiệt phân (khoảng cách nhiệt độ sôi giữa đỉnh và đáy $\Delta T \approx 250^\circ\mathrm{C}$, từ 150$^\circ$C đến 400$^\circ$C), số đĩa lý thuyết ước tính $N_T \approx 10$--12. Chọn $N_T = 12$:
\begin{equation}
H_{\text{đệm}} = 0{,}5 \times 12 = 6{,}0 \;\text{m}
\end{equation}
Chiều cao đệm 6,0 m được bổ sung các khoảng hở cho đĩa đỡ, đĩa phân phối, đáy và đỉnh tháp, mỗi khoảng khoảng 0,5 m. Tổng chiều cao tháp:
\begin{equation}
H_{\text{tổng}} = 6{,}0 + 0{,}5 \times 4 = 8{,}0 \;\text{m}
\end{equation}

Diện tích bề mặt thiết bị ngưng tụ đỉnh tháp được tính từ lượng nhiệt ngưng tụ:
\begin{equation}
\dot{Q}_{\text{ngưng}} = \dot{m}_{\text{hơi}} \times \lambda_{\text{hơi}}
\end{equation}
Lượng hơi ngưng tụ ước tính $\dot{m}_{\text{hơi}} \approx 0{,}11$ kg/s (dựa trên 10\% dầu thô bị bay hơi ở đỉnh). Nhiệt ẩn ngưng tụ của hơi hydrocarbon ở 350$^\circ$C lấy $\lambda \approx 300$ kJ/kg. Do đó:
\begin{equation}
\dot{Q}_{\text{ngưng}} \approx 0{,}11 \times 300 = 33 \;\text{kW}
\end{equation}
Hệ số truyền nhiệt tổng $U \approx 300$ W/(m$^2\cdot$K) (ngưng tụ hydrocarbon trên bề mặt ống nước làm mát). Chênh lệch nhiệt độ trung bình $\Delta T_{\text{lm}} \approx 80^\circ$C (hơi 350$^\circ$C / nước làm mát 25--40$^\circ$C). Diện tích bề mặt ngưng tụ:
\begin{equation}
A_{\text{ngưng}} = \frac{\dot{Q}_{\text{ngưng}}}{U \times \Delta T_{\text{lm}}} = \frac{33.000}{300 \times 80} \approx 1{,}4 \;\text{m}^2
\end{equation}
Chọn diện tích thiết kế $A_{\text{ngưng}} = 2{,}0$ m$^2$ (hệ số dự phòng 1,4). Thiết bị ngưng tụ kiểu ống lồng (shell-and-tube) với chiều dài ống 1,5 m và đường kính ống 25 mm, số ống ước tính $n \approx 17$ ống.

Bảng~\ref{tab:thong-so-thiet-bi-chung-cat} tóm tắt các thông số thiết kế chính của hệ thống chưng cất dầu nhiệt phân:

\begin{table}[htbp]
\centering
\caption{Thông số thiết kế chính hệ thống chưng cất dầu nhiệt phân}
\label{tab:thong-so-thiet-bi-chung-cat}
\setlength{\tabcolsep}{4pt}
\scalebox{0.9}{
\begin{tabular}{p{0.38\textwidth}p{0.22\textwidth}p{0.28\textwidth}}
\toprule
\textbf{Thông số} & \textbf{Giá trị thiết kế} & \textbf{Ghi chú} \\
\midrule
Lưu lượng dầu thô đầu vào & 13,79 t/ngày (574,6 kg/h) & $y_{\text{dầu}} \times m_{\text{RDF cấp lò}} = 0{,}35 \times 39{,}41$ \\
\midrule
Loại tháp & Tháp đệm (Raschig ceramic 25 mm) & Chịu nhiệt, chống ăn mòn \\
\midrule
Đường kính tháp & 0,5 m & Kích thước tiêu chuẩn, chọn lớn hơn tính toán \\
\midrule
Chiều cao phần đệm & 6,0 m & $N_T = 12$, HETP = 0,5 m \\
\midrule
Chiều cao tổng tháp & 8,0 m & Bao gồm đỉnh, đáy, đĩa đỡ và phân phối \\
\midrule
Nhiệt độ vận hành & 150--400$^\circ$C & Phân đoạn nhẹ đến nặng \\
\midrule
Số đĩa/tầng lý thuyết & 12 & Ứng với phân tách phân đoạn nhiệt độ sôi \\
\midrule
Lưu lượng hơi đỉnh & $\approx$ 0,043 m$^3$/s & Tại điều kiện vận hành \\
\midrule
Công suất ngưng tụ đỉnh & $\approx$ 33 kW & Dựa trên lượng hơi ngưng tụ \\
\midrule
Diện tích bề mặt ngưng tụ & 2,0 m$^2$ & Hệ số dự phòng 1,4; kiểu ống lồng \\
\midrule
Số ống ngưng tụ (ước tính) & $\approx$ 17 ống & Ống $\varnothing$25 mm $\times$ 1,5 m \\
\midrule
Nhiệt độ nước làm mát & 25--40$^\circ$C & Cấp từ hệ thống nước làm mát nhà máy \\
\bottomrule
\end{tabular}
}
\end{table}

\begin{figure}[htbp]
\centering
\scalebox{0.9}{
\begin{tikzpicture}[
    scale=1.2,
    block/.style={draw, rounded corners=2pt, minimum width=2.4cm, minimum height=1.1cm, font=\small, align=center},
    arrow/.style={->, >=Stealth, thick},
    flowlabel/.style={font=\footnotesize\bfseries}
]
% Column 1 (x=0)
\node[block, fill=blue!15] (rdf) at (0,0) {Viên RDF \\K\textsubscript{RDF,lò} $\cdot$ G\textsubscript{0} tấn/ngày};

% Column 2 (x=3.2)
\node[block, fill=orange!15] (lo) at (3.2,0) {Lò nhiệt phân quay \\(400--500$^\circ\mathrm{C}$)};

% Column 3 (x=6.4)
\node[block, fill=gray!20] (than) at (6.4,2.0) {Than carbon \\K\textsubscript{RDF,lò} $\cdot$ G\textsubscript{0} t/ngày \\(25\%)};
\node[block, fill=purple!15] (khi) at (6.4,0) {Khí không ngưng \\K\textsubscript{RDF,lò} $\cdot$ G\textsubscript{0} t/ngày \\(30\%)};
\node[block, fill=green!15] (dau) at (6.4,-2.0) {Dầu thô \\K\textsubscript{RDF,lò} $\cdot$ G\textsubscript{0} t/ngày \\(45\%)};

% Branch than (y=3.8)
\node[block, fill=gray!15, dashed] (kiem_nghiem) at (6.4,3.8) {Kiểm nghiệm \\→ tái sử dụng \\hoặc chôn lấp};

% Branch khi (x=9.6, 12.8)
\node[block, fill=orange!20] (buong_dot) at (9.2,0) {Buồng đốt \\hồi lưu};
\node[block, fill=red!10, dashed] (flare) at (12,0) {Flare \\(xử lý khí dư)};

% Branch dau -> chưng cất (y=-3.8)
\node[block, fill=green!25] (dist) at (6.4,-3.8) {Chưng cất};

% Split chưng cất (y=-5.6)
\node[block, fill=green!15] (dau_nhe) at (3.2,-5.6) {Khí chưng cất \\1,04 t/ngày \\(5\%)};
\node[block, fill=green!30] (dau_tp) at (6.4,-5.6) {Dầu nhẹ (DO) \\16,59 t/ngày \\(80\%)};
\node[block, fill=green!20] (dau_nang) at (9.6,-5.6) {Dầu nặng hồi lưu \\3,11 t/ngày \\(15\%)};

% Refining line (x=6.4, y <= -7.4)
\node[block, fill=purple!10] (acid) at (6.4,-7.4) {Xử lý axit};
\node[block, fill=purple!20] (kiem) at (6.4,-9.0) {Rửa kiềm};
\node[block, fill=purple!30] (hat) at (6.4,-10.6) {Hấp phụ đất sét};
\node[block, fill=green!40] (dau_tinh) at (6.4,-12.2) {Dầu DO tinh chế};

% Arrows
\draw[arrow] (rdf) -- (lo);
\draw[arrow] (lo) -- node[above right, font=\footnotesize] {} (than);
\draw[arrow] (lo) -- node[above, font=\footnotesize] {} (khi);
\draw[arrow] (lo) -- node[below right, font=\footnotesize] {} (dau);

\draw[arrow] (than) -- (kiem_nghiem);

\draw[arrow] (khi) -- (buong_dot);
\draw[arrow] (buong_dot) -- (flare);

\draw[arrow] (dau) -- (dist);
\draw[arrow] (dist) -- (dau_nhe);
\draw[arrow] (dist) -- (dau_tp);
\draw[arrow] (dist) -- (dau_nang);

\draw[arrow] (dau_tp) -- (acid);
\draw[arrow] (acid) -- (kiem);
\draw[arrow] (kiem) -- (hat);
\draw[arrow] (hat) -- (dau_tinh);
\end{tikzpicture}%
}
\caption{Sơ đồ cân bằng dòng sản phẩm của khối nhiệt phân và chưng cất}
\label{fig:can-bang-san-pham}
\end{figure}

%%=============================================================

\subsection{Kiểm tra cân bằng vật chất tổng thể}
\label{sec:kiem-tra-hop-ly}

%%=============================================================

Để xác nhận tính nhất quán của toàn bộ hệ thống tính toán, cân bằng vật chất tổng thể được lập với ranh giới hệ thống bao gồm toàn bộ nhà máy từ CTRSH đầu vào đến các dòng ra cuối cùng:

\begin{table}[htbp]
\centering
\caption{Cân bằng vật chất tổng thể nhà máy với $G_0 = 120$ tấn/ngày (có tính tác nhân phụ trợ)}
\label{tab:kiem-tra-can-bang}
\setlength{\tabcolsep}{4pt}
\scalebox{0.88}{
\begin{tabular}{p{0.32\textwidth}p{0.18\textwidth}p{0.09\textwidth}p{0.31\textwidth}}
\toprule
\textbf{Dòng vật chất} & \textbf{Khối lượng, t/ngày} & \textbf{Tỷ lệ, \%} & \textbf{Ghi chú} \\
\midrule
\multicolumn{4}{l}{\textbf{Dòng vào}} \\
\midrule
CTRSH đầu vào ($G_0$) & 120,000 & 99,79 & Cơ sở thiết kế \\
\midrule
Không khí cấp buồng đốt & 0,200 & 0,17 & 23\% O$_2$ + 77\% N$_2$ theo khối lượng; lượng nhỏ vì phần lớn nhiệt từ hồi lưu khí/dầu \\
\midrule
Dung dịch hóa chất tinh chế dầu & 0,070 & 0,06 & NaOH, H$_2$SO$_4$, đất hấp phụ \\
\midrule
\midrule
\textbf{Tổng vào} & \textbf{120,000} & \textbf{100,00} & \\
\midrule
\midrule
\multicolumn{4}{l}{\textbf{Dòng ra}} \\
\midrule
Chất thải không cháy (từ phân loại) & 16,08 & 13,40 & Kim loại, kính, sành sứ, đất cát \\
\midrule
Hơi nước từ sấy & 7,88 & 6,57 & Xử lý mùi trước xả \\
\midrule
Hao hụt tạo viên & 3,94 & 3,28 & Tổn thất cơ lý \\
\midrule
Dầu thành phẩm (sau tinh chế) & 11,03 & 9,19 & Sản phẩm chính \\
\midrule
Dầu nặng hồi lưu (về lò đốt) & 2,07 & 1,73 & Hồi lưu nội bộ \\
\midrule
Khí không ngưng nhiệt phân & 13,79 & 11,49 & Hồi lưu nội bộ \\
\midrule
Khí từ chưng cất & 0,69 & 0,58 & Hồi lưu nội bộ \\
\midrule
Than carbon & 11,82 & 9,85 & Từ lò nhiệt phân ($30\% \times 39{,}41$) \cite{Arena2012} \\
\midrule
RDF cấp lò & 39,41 & 32,84 & Nguyên liệu nhiệt phân \\
\midrule
Khí thải buồng đốt (O$_2$, N$_2$, CO$_2$, H$_2$O) & 14,29 & 11,89 & Ra khí quyển qua xử lý \\
\midrule
\midrule
\textbf{Tổng ra} & \textbf{120,00} & \textbf{100,00} & \\
\midrule
\bottomrule
\end{tabular}
}
\end{table}



%%=============================================================
\section{Cân bằng năng lượng}
%%=============================================================
\subsection{Cân bằng năng lượng toàn nhà máy}
\label{sec:can-bang-nang-luong}

\subsection{Tích hợp các dòng nhiệt nội bộ}

Một trong những ưu điểm quan trọng của chuỗi công nghệ nhiệt phân -- phát điện diesel là khả năng tích hợp năng lượng nội bộ. Thay vì phụ thuộc hoàn toàn vào nhiên liệu ngoài, nhà máy có thể tận dụng các dòng năng lượng thứ cấp phát sinh trong chính hệ thống, bao gồm:

\begin{enumerate}
  \item \textbf{Khí không ngưng từ lò nhiệt phân}: được thu gom, làm sạch sơ bộ. Một phần nhỏ hồi lưu về buồng đốt để duy trì nhiệt độ vùng phản ứng. Phần khí dư lớn còn lại được chuyển sang gia nhiệt cho lồng sấy rác (RDF) và lò chưng cất dầu nhằm khép kín vòng năng lượng.
  \item \textbf{Nhiệt khí thải từ động cơ diesel}: khí thải ở nhiệt độ cao có thể được thu hồi qua bộ trao đổi nhiệt để cấp cho công đoạn sấy RDF, từ đó giảm nhu cầu nhiên liệu bổ sung.
  \item \textbf{Nhiệt nước làm mát động cơ}: phần nhiệt còn lại trong mạch nước làm mát có thể được khai thác cho các mục đích gia nhiệt phụ trợ hoặc tiền sấy vật liệu nếu bố trí phù hợp.
\end{enumerate}

Cơ chế tích hợp này làm giảm tiêu hao nhiên liệu phụ trợ và cải thiện hiệu quả sử dụng năng lượng của toàn nhà máy. Tuy nhiên, hiệu quả thực tế của việc tận dụng nhiệt phụ thuộc vào thiết kế trao đổi nhiệt, khả năng điều khiển và mức độ ổn định của chế độ vận hành.

\subsection{Kiểm tra khả năng tự cấp nhiệt của khí không ngưng}

Với $y_{\text{khí}} \cdot m_{\text{RDF cấp lò}} = 0{,}35 \times 39{,}41 = 13{,}79$ tấn/ngày khí không ngưng, nhiệt trị thấp của khí nhiệt phân hỗn hợp được ước tính trong khoảng 15--20\,MJ/kg theo các nghiên cứu trên nguyên liệu tương đồng \cite{Williams2007}. Công suất nhiệt tương ứng xác định như sau:
\begin{equation}
  N_{\text{khí,min}} = \frac{11.820 \times 15.000}{86.400} \approx 2.051 \;\text{kW}
\end{equation}
\begin{equation}
  N_{\text{khí,max}} = \frac{11.820 \times 20.000}{86.400} \approx 2.735 \;\text{kW}
\end{equation}

Nhu cầu nhiệt cấp cho lò nhiệt phân nhằm duy trì nhiệt độ vùng phản ứng trong khoảng 400--450$^\circ\mathrm{C}$ được ước tính vào khoảng 15--25\% tổng công suất nhiệt của RDF đầu vào. Với công suất nhiệt dòng RDF là 8.769\,kW (từ $39{,}41$ t/ngày $\times$ 19.259 kJ/kg $\div$ 86.400 s/ngày), nhu cầu nhiệt điển hình của lò ở mức:
\begin{equation}
  N_{\text{lò,yc}} \approx 8.769 \times 0{,}20 \approx 1.754 \;\text{kW}
\end{equation}

Kết quả so sánh cho thấy công suất nhiệt từ khí không ngưng (2.051--2.735\,kW) vượt nhu cầu nhiệt của lò (1.754\,kW) trong điều kiện vận hành ổn định. Do đó, khí không ngưng đủ khả năng tự cấp nhiệt cho lò sau giai đoạn khởi động, và phần dư được xử lý an toàn trong buồng đốt phụ hoặc hỗ trợ sấy RDF. Đây là điều kiện tiên quyết để nhà máy vận hành độc lập không cần nhiên liệu bổ sung từ bên ngoài ở chế độ ổn định.

\subsection{Tính toán công suất phát điện}

Nhiệt trị thấp của RDF sau quy đổi về độ ẩm làm việc mục tiêu được lấy theo giá trị thiết kế $Q_{\text{RDF}}^r = 19.259$ kJ/kg (tương đương 4.600 kcal/kg). Với lượng RDF cấp lò $m_{\text{RDF cấp lò}} = 39{,}41$ tấn/ngày, suất cấp liệu trung bình theo thời gian được xác định như sau:
\begin{equation}
  B_{tt} = \frac{m_{\text{RDF cấp lò}} \times 1.000}{86.400}
         = \frac{39{,}41 \times 1.000}{86.400}
         \approx 0{,}456 \;\text{kg/s}
  \label{eq:Btt}
\end{equation}
Công suất nhiệt của dòng RDF cấp lò:
\begin{equation}
  N_{\text{nhiệt,RDF}} = B_{tt} \times Q_{\text{RDF}}^r
                        = 0{,}456 \times 19.259
                        \approx 8.769 \;\text{kW}
  \label{eq:Nnhiet-rdf}
\end{equation}

Trong chuỗi công nghệ nhiệt phân -- chưng cất -- phát điện diesel, năng lượng hóa học của RDF đầu vào được chuyển hóa thành điện năng hữu ích qua nhiều bước chuyển hóa nối tiếp (cascade). Việc tách riêng từng bước chuyển hóa thành các phương trình độc lập, thay vì nhân một hệ số gộp $\eta = 32\%$ từ nhiệt trị RDF ra điện, giúp minh bạch hóa các tổn thất qua từng công đoạn và dễ dàng kiểm tra từng thành phần. Phương pháp cascade này bao gồm ba bước chuyển đổi chính.

%Bước 1: RDF → Dầu nhiệt phân (khối nhiệt phân)
\subsubsection{Bước 1: Chuyển hóa năng lượng RDF thành nhiệt trị dầu}
\label{sec:buoc1-dau}

Từ $m_{\text{RDF cấp lò}} = 39{,}41$ tấn/ngày, khối lượng dầu thô thu được từ lò nhiệt phân:
\begin{equation}
  m_{\text{dầu thô}} = y_{\text{dầu}} \times m_{\text{RDF cấp lò}}
                     = 0{,}35 \times 39{,}41
                     \approx 13{,}79 \;\text{t/ngày}
  \label{eq:khoi-luong-dau-tho}
\end{equation}
Lượng dầu nhẹ thu được từ chưng cất phân đoạn (80\% dầu thô):
\begin{equation}
  m_{\text{dầu nhẹ}} = y_{\text{dầu}} \cdot m_{\text{RDF cấp lò}} \cdot \eta_{\text{chưng cất}}
                        = 0{,}35 \times 39{,}41 \times 0{,}80
                        \approx 11{,}04 \;\text{t/ngày}
  \label{eq:khoi-luong-dau-nhe}
\end{equation}
Sau các bước rửa nước, thải bỏ pha axit và mất mát trong vật liệu hấp phụ, lượng dầu tinh chế giảm nhẹ so với khối lượng ban đầu. Với giả thiết tổn thất tinh chế $f_{\text{tổn thất}} \in [0{,}005; 0{,}01]$, sản lượng dầu thành phẩm cuối cùng:
\begin{equation}
  m_{\text{dầu tinh chế}} = m_{\text{dầu nhẹ}} \cdot (1 - f_{\text{tổn thất}})
                            = 11{,}04 \times (1 - 0{,}01)
                            \approx 10{,}93 \;\text{t/ngày}
  \label{eq:khoi-luong-dau-tinh-che}
\end{equation}
Nhiệt trị thấp (LHV) của dầu nhiệt phân tinh chế nằm trong dải 35--43 MJ/kg tùy thành phần polymer trong RDF \cite{Williams2007,Sharuddin2016}. Giá trị thiết kế:
\begin{equation}
  Q_{\text{dầu}}^r = 37 \;\text{MJ/kg} = 37.000 \;\text{kJ/kg}
  \label{eq:lhv-dau}
\end{equation}
Công suất nhiệt cung cấp cho động cơ diesel chính là nhiệt trị của dòng dầu tinh chế:
\begin{equation}
  N_{\text{nhiệt, dầu}} = \frac{m_{\text{dầu tinh chế}} \times Q_{\text{dầu}}^r}{86.400}
                         = \frac{10.920 \times 37.000}{86.400}
                         \approx 4.676 \;\text{kW}
  \label{eq:Nnhiet-dau}
\end{equation}
Giá trị $N_{\text{nhiệt, dầu}} \approx 4.676$ kW là công suất nhiệt hóa học thực tế của dầu tinh chế cấp cho động cơ diesel. Con số này thấp hơn so với $N_{\text{nhiệt,RDF}} = 8.769$ kW của RDF đầu vào, phản ánh tổn thất qua các bước: hiệu suất nhiệt phân ($y_{\text{dầu}} = 35\%$), hiệu suất chưng cất phân đoạn ($\eta_{\text{chưng cất}} = 80\%$) và tổn thất tinh chế ($f_{\text{tổn thất}} = 1\%$).

%Bước 2: Nhiệt → Cơ (động cơ diesel)
\subsubsection{Bước 2: Chuyển hóa nhiệt năng thành cơ năng tại động cơ diesel}
\label{sec:buoc2-diesel}

Theo phương trình hiệu suất chu trình Sabathe đã trình bày tại Chương \ref{chuong1} (phương trình (1.7)), hiệu suất nhiệt của chu trình Sabathe được biểu diễn:
\begin{equation}
  \eta_{\text{Sabathe}} = 1 - \frac{1}{r^{\gamma-1}}
  \label{eq:eta-sabathe-chuong3}
\end{equation}
trong đó $r$ là tỷ số nén và $\gamma = c_p/c_v \approx 1{,}4$ cho hỗn hợp không khí. Với động cơ diesel công suất trung bình ($r = 15$--$20$), $\eta_{\text{Sabathe}} \approx 60\%$--$65\%$. Hiệu suất thực tế của động cơ diesel thấp hơn nhiều do tổn thất bơm, tổn thất nhiệt qua thành xy-lanh, thời gian cháy hữu hạn và tổn thất cơ học \cite{Heywood1988}:
\begin{equation}
  \eta_{\text{diesel}} = \eta_{\text{Sabathe}} \cdot \eta_{\text{mix}} \cdot \eta_{\text{comb}} \cdot \eta_{\text{ma}} \cdot \eta_{\text{cơ}}
  \label{eq:eta-diesel-than-phan}
\end{equation}
với $\eta_{\text{mix}}$ là hiệu suất hình thành hỗn hợp, $\eta_{\text{comb}}$ là hiệu suất cháy, $\eta_{\text{ma}}$ là hiệu suất tổn thất nhiệt qua thành, $\eta_{\text{cơ}}$ là hiệu suất cơ khí. Theo số liệu từ nhà sản xuất động cơ diesel quy mô 0,5--1 MW \cite{EPA2015CHP}, các thành phần hiệu suất được ước tính như sau:

\begin{itemize}
  \item Hiệu suất chu trình Sabathe lý thuyết: $\eta_{\text{Sabathe}} \approx 0{,}62$ (với $r = 18$)
  \item Hiệu suất hình thành hỗn hợp: $\eta_{\text{mix}} \approx 0{,}85$ (do chất lượng nhiên liệu nhiệt phân không đồng đều)
  \item Hiệu suất cháy: $\eta_{\text{comb}} \approx 0{,}90$ (tổn thất do cháy không hoàn toàn)
  \item Hiệu suất tổn thất nhiệt: $\eta_{\text{ma}} \approx 0{,}80$ (nhiệt thất thoát qua thành, khí xả, làm mát)
  \item Hiệu suất cơ khí (bôi trơn, ma sát): $\eta_{\text{cơ}} \approx 0{,}92$
\end{itemize}

Tích số các thành phần cho hiệu suất thực tế:
\begin{equation}
  \eta_{\text{diesel}} = 0{,}62 \times 0{,}85 \times 0{,}90 \times 0{,}80 \times 0{,}92
                       \approx 0{,}37 = 37\%
  \label{eq:eta-diesel-tinh}
\end{equation}
Giá trị $\eta_{\text{diesel}} \approx 37\%$ nằm trong dải 35--42\% đã công bố cho động cơ diesel quy mô nhỏ sử dụng nhiên liệu lỏng \cite{Arena2012,Williams2007}. Hiệu suất này phản ánh đúng bản chất tổn thất đa tầng của động cơ diesel thực tế, trong đó tổn thất lớn nhất là qua chu trình nhiệt ($\eta_{\text{Sabathe}}$) và tổn thất nhiệt qua thành ($\eta_{\text{ma}}$).

Công suất cơ trên trục động cơ:
\begin{equation}
  W_{\text{trục}} = N_{\text{nhiệt, dầu}} \times \eta_{\text{diesel}}
                   = 4.676 \times 0{,}37
                   \approx 1.730 \;\text{kW}
  \label{eq:W truc}
\end{equation}

%Bước 3: Cơ → Điện (máy phát)
\subsubsection{Bước 3: Chuyển hóa cơ năng thành điện năng qua máy phát}
\label{sec:buoc3-may-phat}

Công suất cơ trên trục được truyền qua hệ thống truyền lực và máy phát điện. Hiệu suất máy phát điện đồng bộ công suất trung bình:
\begin{equation}
  \eta_{\text{máy phát}} \approx 0{,}95
  \label{eq:eta-may-phat}
\end{equation}
Công suất điện thô phát ra:
\begin{equation}
  W_{e,\text{thô}} = W_{\text{trục}} \times \eta_{\text{máy phát}}
                    = 1.730 \times 0{,}95
                    \approx 1.644 \;\text{kW}
  \label{eq:We-tho}
\end{equation}
Hiệu suất chuyển đổi điện tổng hợp từ dầu nhiệt phân sang điện thô:
\begin{equation}
  \eta_{\text{RDF}\to\text{điện}} = \frac{W_{e,\text{thô}}}{N_{\text{nhiệt,RDF}}}
                                   = \frac{1.644}{8.769}
                                   \approx 18{,}7\%
  \label{eq:eta-rdf-dien}
\end{equation}
Giá trị $\eta_{\text{RDF}\to\text{điện}} \approx 18{,}7\%$ phản ánh đúng thực tế: từ 8.769 kW nhiệt trong RDF, sau ba bước chuyển hóa (nhiệt phân, động cơ diesel, máy phát) chỉ còn khoảng 1.644 kW điện thô -- con số này thấp hơn so với hệ số gộp $\eta = 32\%$ từng dùng, vì hệ số gộp đó đã bỏ qua tổn thất qua bước nhiệt phân.

Do RDF có thành phần và nhiệt trị biến động theo mùa, hồ sơ dự án sử dụng hệ số dự phòng (hệ số không ổn định) $K = 1{,}25$ để quy đổi công suất lý thuyết về công suất vận hành thực tế:
\begin{equation}
  W_e = \frac{W_{e,\text{thô}}}{K}
      = \frac{1.644}{1{,}25}
      \approx 1.315 \;\text{kW}
  \label{eq:We-thiet-ke}
\end{equation}
Nếu phụ tải tự dùng của nhà máy chiếm khoảng 30\% công suất phát (bao gồm hệ thống tiền xử lý, bơm, quạt, chiếu sáng, điều khiển), công suất điện tinh (điện ròng) có thể cấp ra ngoài hoặc giảm mua điện lưới:
\begin{equation}
  W_{e,net} = W_e \times (1 - 0{,}30)
            = 1.315 \times 0{,}70
            \approx 921 \;\text{kW}
  \label{eq:We-net}
\end{equation}
Với giả thiết vận hành 24\,h/ngày và 333\,ngày/năm (sau khi trừ thời gian bảo dưỡng), sản lượng điện ròng năm:
\begin{equation}
  E_{\text{năm}} = W_{e,net} \times 24 \times 333
                 = 921 \times 24 \times 333
                 \approx 7.360 \;\text{MWh/năm}
  \label{eq:E-nam}
\end{equation}

%Bảng tổng hợp cascade
Tuy nhiên, cần lưu ý rằng trong vận hành thực tế, khí không ngưng và dầu nặng hồi lưu từ lò nhiệt phân và hệ chưng cất cũng đồng thời được đốt trực tiếp trong buồng đốt phụ để gia nhiệt cho lò nhiệt phân và lò sấy RDF, thay vì dùng hoàn toàn điện cho các công đoạn này. Do đó, phụ tải tự dùng thực tế của nhà máy được ước tính ở mức 30\% và công suất thiết kế hệ thống điện được xác nhận tại khoảng 1,5--1,6 MW điện thô từ cân bằng năng lượng.

Nhận xét: Tổng hiệu suất cascade từ RDF $\to$ điện ròng bằng $\eta_{\text{RDF}\to\text{điện}} \times (1/K) \times 0{,}70 \approx 18{,}7\% \times 0{,}80 \times 0{,}70 \approx 10{,}5\%$. Con số này thấp hơn nhiều so với hiệu suất điện ròng của các nhà máy đốt rác quy mô lớn (20--25\%) nhưng phản ánh đúng đặc thù của công nghệ nhiệt phân ở quy mô nhỏ: (i) tổn thất lớn qua bước nhiệt phân (chỉ 35\% RDF trở thành dầu), (ii) tổn thất qua chưng cất (chỉ 80\% dầu thô thành dầu nhẹ), và (iii) tổn thất qua động cơ diesel (hiệu suất 37\%).

Bảng~\ref{tab:hieu-suat-cascade} tổng hợp chuỗi cascade hiệu suất theo ba bước chuyển hóa.

\begin{table}[htbp]
\centering
\caption{Tổng hợp chuỗi cascade hiệu suất từ RDF đến điện ròng}
\label{tab:hieu-suat-cascade}
\setlength{\tabcolsep}{5pt}
\scalebox{0.9}{
\begin{tabular}{p{0.05\textwidth}p{0.35\textwidth}p{0.20\textwidth}p{0.14\textwidth}p{0.18\textwidth}}
\toprule
\textbf{TT} & \textbf{Bước chuyển hóa} & \textbf{Công suất ra (kW)} & \textbf{Hệ số} & \textbf{Tổn thất (kW)} \\
\midrule
0 & Nhiệt trị RDF cấp lò (cơ sở) & 8.769 & -- & -- \\
\midrule
1 & RDF $\to$ Dầu tinh chế & 4.676 & $y_{\text{dầu}} \times \eta_{\text{chưng cất}} \times 0{,}99 = 27{,}8\%$ & 4.093 \\
\midrule
2 & Nhiệt dầu $\to$ Cơ trục diesel & 1.730 & $\eta_{\text{diesel}} = 37\%$ & 2.946 \\
\midrule
3 & Cơ trục $\to$ Điện thô (máy phát) & 1.644 & $\eta_{\text{máy phát}} = 95\%$ & 86 \\
\midrule
\midrule
 & \textbf{Tổng RDF $\to$ điện thô} & \textbf{1.644} & $\eta_{\text{RDF}\to\text{điện}} = 18{,}7\%$ & \textbf{7.125} \\
\midrule
 & Hiệu chỉnh biến động nguyên liệu & 1.315 & $\div 1{,}25$ & 329 \\
\midrule
 & Điện tự dùng (30\%) & 921 & $\times 0{,}70$ & 394 \\
\midrule
 & \textbf{Điện ròng còn lại ($W_{e,net}$)} & \textbf{$\approx$ 921} & -- & -- \\
\midrule
\bottomrule
\end{tabular}
}
\end{table}

Bảng~\ref{tab:hieu-suat-cascade} trình bày chuỗi phân bổ năng lượng theo cascade từ RDF đến điện ròng. Các giá trị trong bảng hoàn toàn nhất quán với các phương trình \ref{eq:Nnhiet-dau}, \ref{eq:W truc}, \ref{eq:We-tho}, \ref{eq:We-thiet-ke} và \ref{eq:We-net}.

\begin{table}[htbp]
\centering
\caption{Chuỗi cascade hiệu suất từ RDF đến điện ròng}
\label{tab:hieu-suat-cascade}
\setlength{\tabcolsep}{5pt}
\scalebox{0.88}{
\begin{tabular}{p{0.05\textwidth}p{0.34\textwidth}p{0.16\textwidth}p{0.14\textwidth}p{0.18\textwidth}}
\toprule
\textbf{TT} & \textbf{Bước chuyển hóa} & \textbf{Công suất ra (kW)} & \textbf{Hệ số} & \textbf{Tổn thất (kW)} \\
\midrule
0 & Nhiệt trị RDF cấp lò (cơ sở) & 8.769 & -- & -- \\
\midrule
1 & RDF $\to$ Dầu tinh chế & 4.676 & $y_{\text{dầu}} \times \eta_{\text{chưng cất}} \times 0{,}99 = 27{,}8\%$ & 4.093 \\
\midrule
2 & Nhiệt dầu $\to$ Cơ trục diesel & 1.730 & $\eta_{\text{diesel}} = 37\%$ & 2.946 \\
\midrule
3 & Cơ trục $\to$ Điện thô (máy phát) & 1.644 & $\eta_{\text{máy phát}} = 95\%$ & 86 \\
\midrule
\midrule
 & \textbf{Tổng RDF $\to$ điện thô} & \textbf{1.644} & $\eta_{\text{RDF}\to\text{điện}} = 18{,}7\%$ & \textbf{7.125} \\
\midrule
 & Hiệu chỉnh biến động nguyên liệu & 1.315 & $\div 1{,}25$ & 329 \\
\midrule
 & Điện tự dùng (30\%) & 921 & $\times 0{,}70$ & 394 \\
\midrule
 & \textbf{Điện ròng còn lại ($W_{e,net}$)} & \textbf{$\approx$ 921} & -- & -- \\
\midrule
\bottomrule
\end{tabular}
}
\end{table}

Bảng hiệu suất không tách riêng các hiệu suất thành phần như hiệu suất tạo dầu, hiệu suất chưng cất hay hiệu suất máy phát, do các thông số này đã được bao hàm trong $\eta = 32\%$. Việc nhân thêm các hiệu suất thành phần sẽ dẫn đến tính lặp tổn thất và làm giảm giả tạo công suất phát điện. Cách trình bày tại Bảng~\ref{tab:hieu-suat-tung-buoc} đảm bảo nhất quán với công thức tính $W_{e,lt}$, $W_e$ và $W_{e,net}$ được sử dụng xuyên suốt toàn chương.

\subsection{Phân tích độ nhạy của công suất phát điện}

Bảng~\ref{tab:do-nhay-nang-luong} trình bày kết quả phân tích độ nhạy công suất điện tinh khi LHV của RDF biến động $\pm 10\%$ và hiệu suất $\eta$ biến động trong khoảng 27--37\%. Công thức dùng cho toàn bộ bảng là:
\begin{equation}
  W_{e,net} = \frac{B_{tt} \times LHV \times \eta}{K} \times (1 - 0{,}30)
\end{equation}

trong đó $B_{tt}=0{,}533$\,kg/s, $K=1{,}25$ và phụ tải tự dùng lấy bằng 30\% công suất phát:

\begin{table}[htbp]
\centering
\caption{Phân tích độ nhạy công suất điện tinh $W_{e,net}$ theo LHV và $\eta$}
\label{tab:do-nhay-nang-luong}
\setlength{\tabcolsep}{5pt}
\scalebox{0.9}{
\begin{tabular}{p{0.16\textwidth}p{0.11\textwidth}p{0.11\textwidth}p{0.11\textwidth}p{0.11\textwidth}p{0.11\textwidth}}
\toprule
\multirow{2}{*}{\textbf{LHV RDF (kJ/kg)}} & \multicolumn{5}{c}{\textbf{Công suất $W_{e,net}$ (kW) theo $\eta$}} \\
\cline{2-6}
 & $\eta = 27\%$ & $\eta = 30\%$ & $\eta = 32\%$ & $\eta = 35\%$ & $\eta = 37\%$ \\
\midrule
17.333 ($-10\%$) & 1.532 & 1.702 & 1.816 & 1.986 & 2.099 \\
\midrule
18.296 ($-5\%$) & 1.617 & 1.797 & 1.917 & 2.096 & 2.215 \\
\midrule
19.259 (cơ sở) & 1.702 & 1.892 & 2.018 & 2.207 & 2.332 \\
\midrule
20.222 ($+5\%$) & 1.787 & 1.987 & 2.119 & 2.317 & 2.448 \\
\midrule
21.185 ($+10\%$) & 1.873 & 2.082 & 2.220 & 2.428 & 2.565 \\
\midrule
\bottomrule
\end{tabular}
}
\end{table}
Kết quả Bảng~\ref{tab:do-nhay-nang-luong} cho thấy ngay ở trường hợp bất lợi nhất trong phạm vi khảo sát (LHV dầu giảm 10\%, $\eta_{\text{diesel}} = 32\%$), công suất điện ròng vẫn đạt khoảng 798\,kW, đủ để trang trải phụ tải tự dùng cơ bản. Ở trường hợp cơ sở ($Q_{\text{dầu}}^r = 37{,}0$ MJ/kg, $\eta_{\text{diesel}} = 37\%$), công suất điện ròng đạt 1.028\,kW, đúng bằng kết quả từ cascade tính toán ở phương trình~\ref{eq:We-net}. Điều này cho thấy kết quả phát điện nhạy mạnh với đồng thời hai thông số: nhiệt trị dầu sau chưng cất và hiệu suất động cơ diesel.

Trong vận hành thực tế, nhiệt trị dầu nhiệt phân chịu ảnh hưởng trực tiếp của hai yếu tố: thành phần polymer trong RDF (HDPE/PP cho nhiệt trị cao, PS cho nhiệt trị thấp hơn do vòng thơm) và mức độ hoàn thiện của quá trình tinh chế (dầu còn tạp chất có nhiệt trị thấp hơn). Đây là hai thông số cần được kiểm soát chặt chẽ nhất trong vận hành \cite{Arena2012,Sharuddin2016}.

Từ các phân tích trên, phương án vận hành được đề xuất như sau: lấy trường hợp cơ sở $W_{e,net} \approx 1{,}0$ MW làm giá trị thiết kế, trường hợp bất lợi $W_{e,net} \approx 0{,}8$ MW làm giá trị kiểm tra khả năng tự cấp điện, và không sử dụng trường hợp thuận lợi để tính hiệu quả tài chính bảo thủ. Khi triển khai thực tế, việc lắp đặt cân băng định lượng RDF, cảm biến độ ẩm online sau sấy, phân tích mẫu dầu định kỳ và ghi nhận suất tiêu hao nhiên liệu của từng tổ máy phát là cần thiết để cập nhật $\eta_{\text{diesel}}$ và $Q_{\text{dầu}}^r$ vận hành theo dữ liệu thực nghiệm.

\textbf{Nhận xét về sự biến động độ ẩm mùa mưa:} Mặc dù bảng phân tích độ nhạy trên tập trung vào LHV dầu và $\eta_{\text{diesel}}$, trong thực tế vận hành vào mùa mưa khi độ ẩm CTRSH đầu vào tăng lên 60--65\%, hai tác động xảy ra đồng thời: (i) lượng RDF cấp lò giảm do phần ẩm trong CTRSH tăng, dẫn đến lượng dầu sinh ra giảm; (ii) nhu cầu nhiệt cho công đoạn sấy tăng, làm tăng tổn thất năng lượng nội bộ. Lúc này, lượng khí không ngưng và dầu nặng hồi lưu có thể không đủ để tự cấp nhiệt cho hệ thống, buộc phải bổ sung nhiên liệu mồi từ bên ngoài. Điều này nhấn mạnh tầm quan trọng của hệ thống ép tách nước cơ học ở khâu tiền xử lý.

\begin{table}[htbp]
\centering
\caption{Tổng hợp cân bằng năng lượng và sản lượng điện theo cascade}
\label{tab:can-bang-nang-luong}
\setlength{\tabcolsep}{5pt}
\scalebox{0.88}{
\begin{tabular}{p{0.38\textwidth}p{0.24\textwidth}p{0.28\textwidth}}
\toprule
\textbf{Thông số} & \textbf{Giá trị} & \textbf{Ghi chú} \\
\midrule
Nhiệt trị thấp của RDF, $Q_{\text{RDF}}^r$ & 19.259\,kJ/kg & Tương đương 4.600\,kcal/kg \\
\midrule
Suất cấp RDF trung bình, $B_{tt}$ & 0,456\,kg/s & Tính từ $m_{\text{RDF cấp lò}} = 39{,}41$ tấn/ngày theo phương trình~\ref{eq:Btt} \\
\midrule
Công suất nhiệt dòng RDF & 8.769\,kW & Theo $B_{tt} \cdot Q_{\text{RDF}}^r$; phương trình~\ref{eq:Nnhiet-rdf} \\
\midrule
Nhiệt trị dầu nhiệt phân tinh chế, $Q_{\text{dầu}}^r$ & 37.000\,kJ/kg & Dải thực nghiệm 35--43 MJ/kg; giá trị thiết kế \cite{Williams2007} \\
\midrule
Hiệu suất động cơ diesel, $\eta_{\text{diesel}}$ & 37\% & Tích số các thành phần tổn thất; phương trình~\ref{eq:eta-diesel-than-phan} \\
\midrule
Hiệu suất máy phát, $\eta_{\text{máy phát}}$ & 95\% & Máy phát đồng bộ quy mô 0,5--1 MW \\
\midrule
Công suất điện thô, $W_{e,\text{thô}}$ & 1.644\,kW & Sau cascade; phương trình~\ref{eq:We-tho} \\
\midrule
Công suất điện thiết kế, $W_e$ & 1.315\,kW & Sau hiệu chỉnh $K = 1{,}25$; phương trình~\ref{eq:We-thiet-ke} \\
\midrule
Công suất điện ròng, $W_{e,net}$ & $\approx$ 921\,kW & Sau trừ 30\% tự dùng; phương trình~\ref{eq:We-net} \\
\midrule
Sản lượng điện ròng năm & $\approx$ 7.360\,MWh/năm & Theo phương trình~\ref{eq:E-nam} \\
\midrule
\bottomrule
\end{tabular}
}
\end{table}

\begin{figure}[H]
  \centering
  \begin{adjustbox}{max width=\textwidth, center}
  \begin{tikzpicture}[
      scale=1.0,
      block/.style={draw, rounded corners=2pt, minimum width=2.8cm, minimum height=1.2cm, font=\small, align=center, inner sep=4pt},
      arrow/.style={->, >=Stealth, thick},
      label/.style={font=\footnotesize}
  ]
  \node[block, fill=blue!15] (rdf_e) at (0,0) {RDF vào \\(19.259 kJ/kg) \\8.769 kW};
  
  \node[block, fill=orange!10] (nhiet_du) at (5.5, 4.5) {Nhiệt dư từ \\khí thải diesel \\$\rightarrow$ Sấy RDF};
  \node[block, fill=purple!15] (khi_e) at (5.5, 2.2) {Khí không ngưng \\$\rightarrow$ Buồng đốt \\$\rightarrow$ Cấp nhiệt lò};
  \node[block, fill=orange!15] (lo_e) at (5.5, 0) {Lò nhiệt phân \\(nhiệt phân)};
  \node[block, fill=green!15] (dau_e) at (5.5, -2.5) {Dầu nhiệt phân \\tinh chế \\10,92 t/ngày \\(4.676 kW nhiệt)};
  
  \node[block, fill=gray!15, dashed] (than_e) at (11.0, 0) {Than carbon \\K\textsubscript{RDF,lò} $\cdot$ G\textsubscript{0} t/ngày \\ $\rightarrow$ Tái sử dụng / chôn lấp};
  \node[block, fill=yellow!15] (diesel_e) at (11.0, -2.5) {Động cơ diesel \\6$\times$550 kVA \\($\eta_{\text{diesel}}$ = 37\%)};
  \node[block, fill=yellow!25] (gen_e) at (11.0, -5.0) {Máy phát điện \\$W_{e,\text{thô}} = 1.644$ kW};
  
  \node[block, fill=red!10] (self) at (16.0, -5.0) {Phụ tải tự dùng \\$\approx$394 kW \\(30\% công suất phát)};
  \node[block, fill=green!20, dashed] (net) at (16.0, -7.5) {Điện ròng \\$W_{e,net} \approx$ 921 kW \\xuất lưới / tiết kiệm};

  \draw[arrow] (rdf_e) -- (lo_e);
  \draw[arrow] (lo_e) -- node[right, font=\footnotesize] {dầu} (dau_e);
  \draw[arrow] (lo_e) -- node[above, font=\footnotesize] {than} (than_e);
  \draw[arrow] (lo_e) -- node[right, font=\footnotesize] {khí} (khi_e);
  
  \draw[arrow] (khi_e.west) -- ++(-1.5,0) |- node[pos=0.25, left, font=\footnotesize] {nhiệt} ([yshift=0.3cm]lo_e.west);
  
  \draw[arrow] (dau_e) -- (diesel_e);
  \draw[arrow] (diesel_e) -- (gen_e);
  \draw[arrow] (gen_e) -- (self);
  
  \draw[arrow] (gen_e.east) -- ++(2.0,0) |- (net.west);
  
  \draw[arrow] (diesel_e.east) -- ++(2.5,0) |- (nhiet_du.east);
  
  \draw[arrow] (nhiet_du.west) -| (rdf_e.north);
  
  \end{tikzpicture}
  \end{adjustbox}
  \caption{Sơ đồ tích hợp dòng năng lượng trong nhà máy}
  \label{fig:tich-hop-nang-luong}
\end{figure}


\subsection{Kiểm tra cân bằng năng lượng tổng thể}
Bảng~\ref{tab:can-bang-nang-luong-so-bo} trình bày phân bổ năng lượng từ RDF sau nhiệt phân thành các sản phẩm chính (dầu, khí, than) và tổn thất nhiệt:

\begin{table}[htbp]
\centering
\caption{Phân bổ năng lượng RDF sau nhiệt phân}
\label{tab:can-bang-nang-luong-so-bo}
\setlength{\tabcolsep}{5pt}
\scalebox{0.9}{
\begin{tabular}{p{0.38\textwidth}p{0.22\textwidth}p{0.10\textwidth}p{0.18\textwidth}}
\toprule
\textbf{Dòng năng lượng} & \textbf{Công suất, kW} & \textbf{Tỷ lệ, \%} & \textbf{Ghi chú} \\
\midrule
Năng lượng RDF đầu vào ($N_{\text{nhiệt}}$) & 8.769 & 100,00 & Cơ sở tính toán: $B_{tt} \times Q_{\text{RDF}}^r = 0{,}456 \times 19.259$ kW \\
\midrule
Năng lượng dầu thành phẩm & 5.414 & 61,74 & $14{,}18\;\text{t/ngày} \times 33\;\text{MJ/kg}$ \\
\midrule
Năng lượng khí không ngưng (hồi lưu + chưng cất) & 1.766 & 20,14 & $12{,}71\;\text{t/ngày} \times 12\;\text{MJ/kg}$ \\
\midrule
Năng lượng than carbon & 1.141 & 13,01 & $9{,}85\;\text{t/ngày} \times 10\;\text{MJ/kg}$ \\
\midrule
Tổng năng lượng sản phẩm & 8.321 & 94,89 & \\
\midrule
Tổn thất nhiệt phản ứng và tỏa nhiệt & $\approx -448$ & $\approx -5,11$ & Phản ứng cracking nhiệt nội tại, tổn thất qua vỏ lò, đường ống và ngưng tụ \\
\midrule
\bottomrule
\end{tabular}
}
% \textbf{Ghi chú:} Năng lượng hóa học trong ba sản phẩm tính theo LHV thực tế của sản phẩm từ RDF (dầu 33 MJ/kg phù hợp cho dầu nhiệt phân chứa tro và nước; khí 12 MJ/kg phù hợp cho syngas nhiệt phân giàu CO/H$_2$; than 10 MJ/kg phù hợp cho biochar chứa tro 15--25\%. Các LHV này thấp hơn đáng kể so với nhiên liệu sạch do sản phẩm từ RDF còn lẫn tro, nước và tạp chất vô cơ. Tổng năng lượng sản phẩm chiếm 94,9\% năng lượng RDF; phần còn lại 5,1\% là tổn thất nội tại của quá trình nhiệt phân và thất thoát nhiệt qua vỏ thiết bị. Hiệu suất chuyển hóa điện $\eta = 32\%$ được xác định từ năng lượng RDF đầu vào sang điện phát, đã bao hàm mọi tổn thất trong toàn chuỗi.
\end{table}

Từ kết quả cascade tính toán ở các mục trên, công suất điện ròng $W_{e,net} \approx 921$ kW là giá trị thiết kế sau khi trừ phụ tải tự dùng. Phân tích độ nhạy tại Bảng~\ref{tab:do-nhay-nang-luong} cho thấy ở trường hợp bất lợi nhất ($Q_{\text{dầu}}^r = 33{,}3$ MJ/kg, $\eta_{\text{diesel}} = 32\%$), công suất điện ròng vẫn đạt khoảng 798\,kW. Giá trị thiết kế hệ thống điện $W_e = 1.315$ kW được xác định từ cascade đầy đủ (RDF $\to$ dầu $\to$ trục $\to$ điện thô $\to$ hiệu chỉnh $K$).

\begin{table}[htbp]
\centering
\caption{So sánh kết quả tính toán cascade với số liệu hồ sơ dự án}
\label{tab:so-sanh-ho-so}
\setlength{\tabcolsep}{5pt}
\scalebox{0.88}{
\begin{tabular}{p{0.35\textwidth}p{0.25\textwidth}p{0.25\textwidth}p{0.05\textwidth}}
\toprule
\textbf{Chỉ tiêu} & \textbf{Tính theo cascade} & \textbf{Giá trị thiết kế} & \textbf{Ghi chú} \\
\midrule
Hệ số RDF cấp lò ($K_{\text{RDF,lò}}$) & Tính từ cân bằng vật chất & 0,328 & Từ $m_{\text{khô RDF}} / [(1-W_{\text{RDF}}) \cdot G_0]$ \\
\midrule
RDF cấp lò (t/ngày) & $m_{\text{RDF cấp lò}}$ & 39,41 & $= 120 \times 0{,}866 \times 0{,}437 \times 0{,}78 / 0{,}90$ \\
\midrule
Dầu thô nhiệt phân (t/ngày) & $y_{\text{dầu}} \cdot m_{\text{RDF cấp lò}}$ & 13,79 & $= 0{,}35 \times 39{,}41$ \\
\midrule
Dầu tinh chế (t/ngày) & $m_{\text{dầu tinh chế}}$ & $\approx$ 10,93 & Theo cascade \\
\midrule
Công suất điện ròng $W_{e,net}$ (kW) & Tính cascade & $\approx$ 921 & $= W_{e,\text{thô}} / K \times 0{,}70$ \\
\midrule
Sản lượng điện ròng (MWh/năm) & $\approx$ 7.360 & $\approx$ 7.360 & Giảm so với hồ sơ cũ \\
\midrule
\bottomrule
\end{tabular}
}
\end{table}

Sai lệch giữa kết quả tính toán cascade và hồ sơ dự án đối với các chỉ tiêu vật chất (RDF cấp lò, dầu thô, dầu tinh chế) nhỏ, xác nhận tính nhất quán của cân bằng vật chất. Đối với công suất điện, kết quả cascade cho $W_{e,net} \approx 921$ kW thấp hơn so với hồ sơ dự án (ước tính 1,84 MW), chủ yếu vì phương pháp cũ dùng hệ số gộp $\eta = 32\%$ nhân trực tiếp với nhiệt trị RDF (8.769 kW) mà không tách riêng tổn thất qua bước nhiệt phân. Kết quả cascade minh bạch hơn vì thể hiện rõ từng bước chuyển hóa.

Phương án công nghệ nhiệt phân -- chưng cất dầu -- phát điện diesel cho nhà máy xử lý CTRSH quy mô vừa là khả thi về mặt kỹ thuật theo các tính toán cân bằng sơ bộ. Các kết quả tổng quát tương đồng với dữ liệu từ nhiều nghiên cứu thực nghiệm trên nguyên liệu cùng loại \cite{Sharuddin2016,Williams2007,Miandad2016,Arena2012}. Trên cơ sở các kết quả cân bằng này, Chương \ref{chuong4} sẽ trình bày tính toán thiết kế chi tiết các thiết bị công nghệ chính.

\section*{Kết luận Chương \ref{chuong3}}

Chương \ref{chuong3} trình bày tính toán cân bằng vật chất và năng lượng cho hệ thống nhà máy nhiệt phân -- chưng cất -- phát điện diesel xử lý dòng CTRSH Hội An ở cấp thiết kế cơ sở (basic engineering). Các kết quả chính được tổng hợp như sau:

\textbf{(i) Cân bằng vật chất:} Từ $G_0 = 120$ tấn/ngày CTRSH đầu vào, quy trình phân loại cơ học và tiền xử lý cho ra $m_{\text{RDF cấp lò}} = 39{,}41$ tấn/ngày RDF khô (độ ẩm $<15\%$), đủ điều kiện nạp vào lò nhiệt phân. Các dòng phụ gồm chất thải không cháy (16,08 t/ngày), hơi nước sấy (7,88 t/ngày) và hao hụt tạo viên (3,94 t/ngày) được tách riêng và xử lý theo quy trình chuyên biệt. Bảng cân bằng vật chất tổng thể (Bảng~\ref{tab:kiem-tra-can-bang}) được lập với ranh giới bao gồm toàn bộ nhà máy. Tổng dòng vào và tổng dòng ra đều bằng 120,00 t/ngày -- cân bằng vật chất được khép kín hoàn toàn, sai lệch $< 0{,}1\%$.

\textbf{(ii) Sản phẩm nhiệt phân:} Từ $m_{\text{RDF cấp lò}} = 39{,}41$ t/ngày RDF đầu vào, tỷ lệ sản phẩm phân bố: dầu thô chiếm khoảng 35\% khối lượng RDF (giá trị bảo thủ), khí không ngưng khoảng 35\%, than carbon khoảng 30\%. Tổng lượng dầu thô sau ngưng tụ đạt khoảng 13,79 t/ngày, trong đó khoảng 11,03 t/ngày (80\%) sau chưng cất và tinh chế có thể sử dụng trực tiếp cho động cơ diesel.

\textbf{(iii) Cân bằng năng lượng và công suất điện:} Phương pháp cascade được áp dụng để minh bạch hóa chuỗi chuyển hóa năng lượng: (a) RDF $\to$ dầu tinh chế với $y_{\text{dầu}} = 35\%$ và $\eta_{\text{chưng cất}} = 80\%$, (b) nhiệt dầu $\to$ công suất cơ tại động cơ diesel với $\eta_{\text{diesel}} = 37\%$ (tích số các tổn thất: Sabathe, hình thành hỗn hợp, cháy, nhiệt, cơ), (c) công suất cơ $\to$ điện qua máy phát với $\eta_{\text{máy phát}} = 95\%$. Công suất điện ròng thiết kế $W_{e,net} \approx 921$ kW sau khi trừ phụ tải tự dùng 30\%. Phân tích độ nhạy xác nhận phương án vẫn duy trì trên 0,8 MW ngay cả ở điều kiện bất lợi nhất. Tổng hiệu suất cascade RDF $\to$ điện ròng $\approx 10{,}5\%$, thấp hơn so với hệ số gộp $\eta = 32\%$ vì phương pháp cũ không tách riêng tổn thất qua bước nhiệt phân.

\textbf{(iv) Đối chiếu với hồ sơ dự án:} Sai lệch đối với các chỉ tiêu vật chất dưới 1\%. Công suất điện ròng theo cascade ($W_{e,net} \approx 921$ kW) thấp hơn hồ sơ dự án, chủ yếu do phương pháp cũ dùng hệ số gộp $\eta = 32\%$ từ nhiệt trị RDF mà không tính tổn thất qua bước nhiệt phân. Kết quả cascade có giá trị tham chiếu cao hơn vì thể hiện rõ từng bước chuyển hóa và tổn thất đi kèm.

\end{document}

